\documentclass[11pt]{ctexart}
\usepackage[a4paper,margin=1in]{geometry}
\usepackage{graphicx}
\usepackage{xcolor}
\usepackage{hyperref}
\definecolor{refgray}{gray}{0.45}
\title{\textbf{Market Views｜全球投行叙事汇编}\\[0.4em]{\large AI资本开支从云端向存储、设备与电力全链条传导，美元承压与新兴市场分化并行}}
\author{KC桌面 自动生成}
\date{260801}
\begin{document}
\maketitle
\tableofcontents
\newpage
\section{一页摘要}
\begin{itemize}
  \item AI资本开支周期仍在扩张，但增长逻辑正从单一GPU转向多元算力与系统级能力，云厂商资本开支维持高位，供给仍跟不上需求。
  \item 存储行业由价格周期转向订单周期，长期协议锁定产能，供需缺口或延续至2028年，HBM4成为竞争焦点。
  \item AI服务器需求收紧MLCC与高端CCL供需，订单能见度拉长至2027年，高端产品占比提升推动毛利率结构性上行。
  \item AI数据中心电力需求将燃气轮机、气体发电机等设备订单排至2030年代，电网建设进入兑现期，承包商和设备商上调指引。
  \item 美元短期承压但利差保护提供缓冲，美联储信誉不确定性导致曲线扭曲陡峭化，新兴市场货币在油价与利差交易中精细化分化。
  \item 中国出口受AI相关产品价格驱动显著，价格贡献远大于数量，官方PMI与出口导向PMI背离反映内外需动能差异。
  \item 高端品牌定价权稳固，大众消费在中国市场承压，啤酒与运动服饰分化，产品结构升级提供局部亮点。
  \item 跨国药企对中国创新药资产投入从许可引进转向大规模后期临床开发，CRO受益于新兴生物制药融资回暖。
\end{itemize}
\section{全球投行叙事汇编}
今日纳入 92 篇投行/券商研究，按 20 个市场、资产与产业主题系统整理。
\newpage
\subsection{宏观与市场策略：AI驱动、美元承压与新兴市场分化}
\textbf{全球宏观主线正从总量增长转向结构性分化：AI资本开支与贸易重塑中国出口和半导体设备需求，美联储信誉问题引发美元曲线扭曲陡峭化，新兴市场货币在油价与利差交易中呈现精细化分化。}
\subsubsection*{主流共识}
\begin{itemize}
  \item AI资本开支周期仍在扩张，微软等超大规模平台维持高投入，支撑全球科技供应链。
  \item 中国出口受AI相关产品价格驱动显著，价格贡献远大于数量。
  \item 美元短期承压，但利差保护与加息预期提供缓冲，趋势未反转。
  \item 亚洲保险业杠杆整体可控，但资本效率与资本充足存在权衡。
\end{itemize}
\subsubsection*{分歧与条件差异}
\begin{itemize}
  \item 中国官方PMI与出口导向PMI背离，反映内需与外需动能差异。
  \item 新兴市场货币对油价冲击的反应分化，能源出口国与进口国表现不一。
  \item 美联储信誉不确定性导致曲线扭曲陡峭化，但短端利率重定价有限。
  \item AI设备增长引擎从存储转向逻辑，但C2W放量时间点存在不确定性。
\end{itemize}
\subsubsection*{机构观点}
\paragraph{BARC} BARC将开云集团评级从低配上调至市场权重，核心逻辑是经营企稳与去杠杆超预期。新CEO上任后迅速出售Kering Beauté回笼40亿欧元，净债务从80亿降至33亿欧元，净杠杆从3.4倍降至1.4倍。Gucci Q2有机降幅收窄至2\%，好于预期，皮具恢复增长，经营利润率同比提升100bp至17\%。但BARC同时将欧洲高等级零售板块下调至低配，因亚马逊进入该市场改变板块格局。
{\small\color{refgray}数据：净债务从80亿欧元降至33亿欧元\par}
{\small\color{refgray}数据：净杠杆从3.4倍降至1.4倍\par}
{\small\color{refgray}数据：Gucci Q2有机收入降2\%，好于预期的4.7\%\par}
{\small\color{refgray}数据：Gucci经营利润率17\%，同比+100bp\par}
{\small\color{refgray}边际变化：从低配上调至市场权重，去杠杆速度快于预期，经营出现企稳信号。\par}
\paragraph{JPM} JPM上调ASMPT 2026-2028年EPS预测7\%-12\%，核心是逻辑TCB设备订单放量、经营杠杆改善和光电子业务贡献。2Q26收入49.36亿港元，环比+24\%，经营利润率15.9\%远超预期。7月单月来自头部OSAT客户的C2S逻辑TCB订单超50台，市占率近100\%，支撑三季度订单环比高个位数增长。C2W有望在2027年下半年成为下一增长点，HBM4延迟后下半年订单重启。
{\small\color{refgray}数据：2Q26收入49.36亿港元，环比+24\%\par}
{\small\color{refgray}数据：经营利润率15.9\%，远超预期的10.2\%\par}
{\small\color{refgray}数据：7月C2S TCB订单超50台\par}
{\small\color{refgray}数据：2026-28年EPS上调7\%-12\%\par}
{\small\color{refgray}边际变化：增长叙事从HBM4转向逻辑先进封装和光电子，C2W成为新关注点。\par}
\paragraph{JPM} JPM亚洲保险报告将杠杆拆分为承保、资产、财务三维度，认为除非严重宏观冲击，整体杠杆可控。亚洲保险公司更关注资本充足而非效率，承保杠杆保守源于低利润储蓄型产品和监管约束。资产杠杆方面，亚洲寿险平均14倍、非寿险4倍，印度偏高。准备金计提保守的公司（AIA、三星生命、Thai Life）享有估值溢价。
{\small\color{refgray}数据：亚洲寿险资产杠杆平均14倍\par}
{\small\color{refgray}数据：非寿险4倍\par}
{\small\color{refgray}数据：承保利润率通常低于10\%\par}
{\small\color{refgray}数据：ROE可达13\%左右\par}
{\small\color{refgray}边际变化：强调杠杆质量（准备金历史）而非单纯杠杆水平，区分估值分化。\par}
\paragraph{JPM} JPM外汇团队在FOMC后维持做多美元，但承认美联储信誉问题引入通胀不确定性溢价。2s30s曲线单日扭曲陡峭化15bp，为十年来FOMC日最大，盈亏平衡通胀率跳升4-7bp，美元在发布会期间走弱。但短端OIS仅下行5-6bp，12月加息预期和地缘贸易条件支撑美元。维持做多美元兑G10低息货币篮子。
{\small\color{refgray}数据：2s30s单日走阔15bp\par}
{\small\color{refgray}数据：盈亏平衡通胀率+4-7bp\par}
{\small\color{refgray}数据：1年内OIS下行5-6bp\par}
{\small\color{refgray}数据：12月加息预期升温\par}
{\small\color{refgray}边际变化：美元回调被视为信用折价而非趋势转变，关注后续非农数据。\par}
\paragraph{NOM} NOM预计7月中国官方制造业PMI回落至49.9，而出口导向的RatingDog PMI维持51.5，反映内外需分化。工业增加值增速预计放缓至5.0\%，出口增速22.9\%仍高，进口31.5\%受芯片价格推动。政治局会议强调稳增长，但政策规模可能有限。AI繁荣与房地产K型分化，一线城市新房回暖而低线城市跌幅扩大。
{\small\color{refgray}数据：官方PMI预计49.9\par}
{\small\color{refgray}数据：RatingDog PMI预计51.5\par}
{\small\color{refgray}数据：出口增速预计22.9\%\par}
{\small\color{refgray}数据：进口增速预计31.5\%\par}
{\small\color{refgray}边际变化：官方与出口PMI背离持续，政策刺激规模受限。\par}
\paragraph{NOM} NOM认为微软4QFY26财报缓解AI基础设施资本开支可持续性担忧。Azure收入增长43\%，云收入593亿美元，下季度capex指引超500亿美元，FY27E继续增长且自由现金流保持正值。RPO剔除OpenAI后增长25\%，显示AI需求客户基础扩散。相关标的中际旭创、生益科技、沪电股份受益。
{\small\color{refgray}数据：Azure收入+43\%\par}
{\small\color{refgray}数据：云收入593亿美元，+27\%\par}
{\small\color{refgray}数据：下季度capex指引500亿+\par}
{\small\color{refgray}数据：RPO剔除OpenAI后+25\%\par}
{\small\color{refgray}边际变化：资本开支指引不降反升，回应市场对AI投入可持续性的疑虑。\par}
\paragraph{NOM} NOM中间价模型预测美元兑人民币中间价6.7300，较前值下调592个基点，含逆周期因子为6.7426。调整幅度显著，反映央行对近期市场供求和外部因素响应加快。逆周期因子修正126个基点，显示央行对汇率波动的容忍区间。关注8月货币政策报告、9月美联储会议等节点。
{\small\color{refgray}数据：预测中间价6.7300\par}
{\small\color{refgray}数据：较前值下调592bp\par}
{\small\color{refgray}数据：含逆周期因子6.7426\par}
{\small\color{refgray}数据：逆周期因子修正126bp\par}
{\small\color{refgray}边际变化：单次预测调整幅度明显，中间价定价机制响应速度加快。\par}
\paragraph{UBS} UBS指出2026年上半年中国出口增长17.6\%、进口增长26\%，AI相关产品贡献约一半增量，但主要由价格驱动。电子集成电路出口金额+96\%而数量仅+6\%，反映供给受限下的价格效应。AI相关贸易超70\%在亚洲内部，香港为转口枢纽。非AI出口仍增长11\%，汽车出口+60\%为亮点。
{\small\color{refgray}数据：出口+17.6\%，进口+26\%\par}
{\small\color{refgray}数据：AI贡献出口增量49\%、进口47\%\par}
{\small\color{refgray}数据：集成电路出口金额+96\%，数量+6\%\par}
{\small\color{refgray}数据：非AI出口+11\%，汽车+60\%\par}
{\small\color{refgray}边际变化：AI贸易成为核心驱动，价格因素主导，深度嵌入亚洲供应链。\par}
\paragraph{GS} GS分析7月油价见顶后新兴市场货币表现，本轮反应幅度小于3月，因不确定性偏好未同步恶化、多数进口国货币已处年内低位、贸易条件冲击较小。模型显示多数货币实际表现略好于预测，韩元和哥伦比亚比索正向偏离，南非兰特、智利比索和墨西哥比索负向偏离。利差交易仍在运行，但需精细选择。
{\small\color{refgray}数据：新兴市场指数从6月高点回撤12\%\par}
{\small\color{refgray}数据：剔除科技后7月涨2\%\par}
{\small\color{refgray}数据：韩元受益于政策收紧\par}
{\small\color{refgray}数据：南非兰特因央行按兵不动跑输\par}
{\small\color{refgray}边际变化：汇率表现与模型偏离度成为识别本地因素的关键框架。\par}
\subsubsection*{关键数据}
\begin{itemize}
  \item 开云净债务从80亿降至33亿欧元，净杠杆1.4倍
  \item ASMPT 7月C2S TCB订单超50台，市占率近100\%
  \item 2s30s FOMC日扭曲陡峭化15bp，十年最大
  \item 中国官方PMI预计49.9 vs 出口PMI 51.5
  \item 微软下季度capex指引超500亿美元
  \item 人民币中间价预测下调592bp至6.7300
  \item AI贡献中国出口增量49\%，集成电路金额+96\%但数量+6\%
  \item 新兴市场指数从高点回撤12\%，剔除科技后7月涨2\%
\end{itemize}
\begin{figure}[htbp]
\centering
\includegraphics[width=0.88\linewidth]{figures/fig_066.jpg}
\caption{Figure 1 - Figure 1: USD weakness matched the contours of 2s30s twist steepening LHS: 2s30s; RHS: DXY (inverted) \#\# Global FX Strategy Arindam Sandilya JPM Chase Bank, N.A., Singapore Branch}
\end{figure}
\begin{figure}[htbp]
\centering
\includegraphics[width=0.88\linewidth]{figures/fig_051.jpg}
\caption{Figure 1 - Figure 1: ASMPT annual revenue and yoy trend Figure 2: ASMPT quarterly revenue and yoy trend Figure 3: ASMPT gross profit and GM trend}
\end{figure}
\begin{figure}[htbp]
\centering
\includegraphics[width=0.88\linewidth]{figures/fig_102.jpg}
\caption{Figure 1 - Figure 1: Price increase explains a large portion of trade strength YTD Figure 2: RMB strengthened in EER terms \#\# AI trade has emerged as a key driver}
\end{figure}
\begin{figure}[htbp]
\centering
\includegraphics[width=0.88\linewidth]{figures/fig_146.jpg}
\caption{Exhibit 1 - Exhibit 2: Many energy importers were trading close to their weakest levels YTD even before the re-escalation Broader terms of trade shift also smaller than in early March, with many energy importers' indices above their weakest}
\end{figure}
{\small\color{refgray}本节综合 10 篇研究；涉及 BARC、JPM、NOM、UBS、GS。}
\newpage
\subsection{AI算力需求的结构性扩张与供应链传导}
\textbf{AI基础设施投入仍在加速，但增长逻辑正从单一GPU转向多元算力与系统级能力，供应链各环节的景气度与风险点出现分化。}
\subsubsection*{主流共识}
\begin{itemize}
  \item AI算力需求持续旺盛，云厂商资本开支维持高位，供给仍跟不上需求。
  \item AWS等云业务增速加快，AI相关收入成为核心增长引擎。
  \item AI芯片测试需求从GPU扩散至CPU和ASIC，测试设备行业景气度上行。
  \item 国产AI芯片市场长期增长确定，但短期交付执行是主要约束。
\end{itemize}
\subsubsection*{分歧与条件差异}
\begin{itemize}
  \item AI资本开支的可持续性：Meta上调指引但幅度有限，微软因会计调整下调数字但实际计划未变，市场对过度投入的担忧仍存。
  \item AI变现路径：Meta强调出售Intelligence利润率高于Compute，而AWS以算力租赁为主，两种模式对产业链的拉动结构不同。
  \item 测试设备格局：Teradyne在二供认证上取得进展，但份额迁移最早2027年才显现，短期格局仍由Advantest主导。
  \item Arm版税增速放缓与AGI CPU供应不确定性，市场对Arm AI故事的分歧加大。
\end{itemize}
\subsubsection*{机构观点}
\paragraph{Bernstein} AWS剔除AI后的非AI收入环比增速约6\%，与Datadog核心客户模型高度吻合（相关性0.9倍以上）。据此推算，Datadog本季度核心收入增速将从约26.0\%升至27\%低至中段。AI收入占比提升削弱了AWS整体增速对Datadog的指向性，但非AI部分仍是有效领先指标。Q4面临基数变化，同比增速预计有近200个基点的波动。
{\small\color{refgray}数据：AWS二季度收入同比增36.7\%，AI年化收入从150亿增至250亿美元\par}
{\small\color{refgray}数据：AWS非AI收入环比增速约6\%\par}
{\small\color{refgray}数据：Datadog核心收入增速预计从26.0\%升至27\%低至中段\par}
{\small\color{refgray}数据：Q4同比增速预计有近200个基点波动\par}
{\small\color{refgray}边际变化：强调剔除AI后的AWS增速才是核心信号，而非整体高增长。\par}
\paragraph{JPM} AI推理需求正同时拉动ASIC、CPU和DRAM三类芯片的测试需求，增速快于预期。爱德万上调2026年SoC测试机市场预期并加速扩产，客户订单可见度从6个月延长至18个月。测试机与治具在FT环节1:1配套，鸿劲精密作为AI/HPC FT治具龙头直接受益。2025年AI GPU与ASIC需求比例约75:25，2026年预计变为50:50，CPU出货量增长显著。
{\small\color{refgray}数据：爱德万V93K产能目标提前完成（5000台原定2027年3月）\par}
{\small\color{refgray}数据：客户订单可见度从6个月延长至18个月\par}
{\small\color{refgray}数据：2025年AI GPU:ASIC需求比例75:25，2026年预计50:50\par}
{\small\color{refgray}数据：FT环节测试机与治具1:1配套\par}
{\small\color{refgray}边际变化：AI测试需求从GPU独大转向CPU/ASIC多元化，需求可见度大幅拉长。\par}
\paragraph{JPM} Teradyne在超大规模客户二供认证上取得进展，已向一家GPU客户交付多系统测试机，并正争取某大型云厂商的第二供应商认证（第一家为Advantest）。但份额迁移将是渐进的，认证周期约9-12个月，明显变化最早2027年才会出现。Teradyne对2026年ATE市场判断为同比增长20-40\%，SoC测试机增速高于存储器。
{\small\color{refgray}数据：Teradyne 4-6月销售额13.29亿美元，超一致预期12.17亿\par}
{\small\color{refgray}数据：7-9月指引12-13亿美元，高于一致预期10.35亿\par}
{\small\color{refgray}数据：2026年ATE市场预计同比增长20-40\%\par}
{\small\color{refgray}数据：GPU与CPU测试强度比约4:1\par}
{\small\color{refgray}边际变化：SoC测试机从单供应商向双供应商模式转变的早期迹象。\par}
\paragraph{MS} Arm版税增速放缓，下季度指引约13\%低于预期，但数据中心版税同比增长超100\%。AGI CPU销售预期上调至FY27-28超10亿美元，需求管道超20亿美元，但供应覆盖是当前主要约束，需到Q3才能确认。当前报价已较大程度反映新CPU业务机会。
{\small\color{refgray}数据：下季度版税指引约13\%\par}
{\small\color{refgray}数据：数据中心版税同比增长超100\%\par}
{\small\color{refgray}数据：AGI CPU FY27-28销售预期超10亿美元，需求管道超20亿美元\par}
{\small\color{refgray}数据：供应覆盖需Q3确认\par}
{\small\color{refgray}边际变化：版税增速放缓，AGI CPU从需求故事转向供应约束。\par}
\paragraph{MS} 寒武纪MLU580/590交付进度慢于预期，供应链检查显示生产爬坡、存储供应和系统级交付不成熟。预计2Q26收入38亿元（环比+31\%），全年收入预测下调2\%，但毛利率假设维持（2026-2028年51\%/50\%/49\%）。股权激励计划考核目标低于市场预期，更接近留人机制而非业绩指引。国产AI芯片竞争正从单芯片转向系统级能力。
{\small\color{refgray}数据：2Q26收入预计38亿元，环比+31\%\par}
{\small\color{refgray}数据：全年收入预测下调2\%\par}
{\small\color{refgray}数据：2030年中国AI芯片TAM上调至910亿美元（CAGR 23\%）\par}
{\small\color{refgray}数据：股权激励覆盖85\%员工，授予价750元/股\par}
{\small\color{refgray}边际变化：交付延迟是短期执行问题，长期需求逻辑未变。\par}
\paragraph{NOM} Meta 2Q26营收608亿美元（同比+28\%），但净利润同比-13.6\%，主因资本开支加速（单季301亿美元，同比+82\%）。全年资本开支指引收窄至1300-1450亿美元。管理层明确表示出售Intelligence的利润率将持续显著高于出售Compute，AI变现路径从卖算力转向卖智能。
{\small\color{refgray}数据：Meta 2Q26资本开支301亿美元，同比+82\%\par}
{\small\color{refgray}数据：全年资本开支指引1300-1450亿美元\par}
{\small\color{refgray}数据：经营利润率30.9\%，低于预期4.6个百分点\par}
{\small\color{refgray}数据：Advantage+年化收入突破750亿美元\par}
{\small\color{refgray}边际变化：AI变现叙事从算力转向智能，利润率差异成为关键。\par}
\paragraph{NOM} AWS 2Q26收入422亿美元（同比+36.8\%），积压订单4960亿美元（同比三位数增长），年化收入1690亿美元。全年资本开支指引上调至2200亿美元（内存成本上升），2028年部分产能已被预订，服务器和网络设备回收期不到三年。芯片业务年化收入超250亿美元，Trainium客户范围扩大。
{\small\color{refgray}数据：AWS 2Q26收入422亿美元，同比+36.8\%\par}
{\small\color{refgray}数据：积压订单4960亿美元，同比三位数增长\par}
{\small\color{refgray}数据：全年资本开支指引2200亿美元\par}
{\small\color{refgray}数据：芯片业务年化收入超250亿美元\par}
{\small\color{refgray}边际变化：云产能需求可见度延伸至2028年，缓解过度投入担忧。\par}
\paragraph{NOM} Meta和微软均表示维持混合模型策略，不局限于特定模型，开源权重模型与自研前沿模型并行。两家公司均称供给跟不上需求，但资本开支未超预期。Meta上调全年指引至1300-1450亿美元，微软因会计调整从1900亿下调至1750亿但实际计划未变。
{\small\color{refgray}数据：Meta 2Q26资本开支301亿美元，低于预期331亿\par}
{\small\color{refgray}数据：微软4Q26资本开支358亿美元\par}
{\small\color{refgray}数据：Meta全年指引1300-1450亿美元\par}
{\small\color{refgray}数据：微软全年指引1750亿美元（会计调整）\par}
{\small\color{refgray}边际变化：模型策略从单一押注转向多模型并行，算力需求结构更复杂。\par}
\paragraph{UBS} AWS AI年化收入从一季度100亿跳升至250亿美元，backlog达4960亿美元（剔除Anthropic 100亿协议后仍超预期510亿）。据此上调AWS增速预测：3Q26从35.8\%上调至38.9\%，全年从33.3\%上调至36.1\%，2027年将突破40\%（OpenAI 2027年初使用Trainium）。剔除能源成本收益后AWS经营利润率持平在38\%。
{\small\color{refgray}数据：AWS AI年化收入250亿美元（上季度100亿）\par}
{\small\color{refgray}数据：backlog 4960亿美元，剔除Anthropic后超预期510亿\par}
{\small\color{refgray}数据：AWS 3Q26增速预测上调至38.9\%\par}
{\small\color{refgray}数据：2027年增速预计超40\%\par}
{\small\color{refgray}边际变化：AWS增长叙事从AI概念兑现推进到核心云与AI算力需求同步加速。\par}
\subsubsection*{关键数据}
\begin{itemize}
  \item AWS AI年化收入从100亿增至250亿美元（R051）
  \item AWS积压订单4960亿美元，同比三位数增长（R046/R051）
  \item Meta 2Q26资本开支301亿美元，同比+82\%（R045）
  \item 爱德万客户订单可见度从6个月延长至18个月（R019）
  \item Teradyne 2026年ATE市场预计同比增长20-40\%（R024）
  \item Arm数据中心版税同比增长超100\%（R033）
  \item 寒武纪2Q26收入预计38亿元，环比+31\%（R036）
  \item 2030年中国AI芯片TAM上调至910亿美元（R036）
\end{itemize}
\begin{figure}[htbp]
\centering
\includegraphics[width=0.88\linewidth]{figures/fig_040.jpg}
\caption{Exhibit 1 - AWS web metrics update: July starting behind, but good trajectory (in line with our model expectation). A quick update on AWS web metrics that we use on a +1 quarter basis to preview how we think core growth could evolve. As we've covered before, Q2 growth fla}
\end{figure}
\begin{figure}[htbp]
\centering
\includegraphics[width=0.88\linewidth]{figures/fig_076.jpg}
\caption{Figure 1 - Figure 2: Advantest and Teradyne's QoQ sales trends (CY basis) Note: CY basis. The definition of tester sales differs between the two companies, but figures are based on disclosed data.}
\end{figure}
\begin{figure}[htbp]
\centering
\includegraphics[width=0.88\linewidth]{figures/fig_084.jpg}
\caption{Exhibit 3 - Exhibit 3: We expect China AI chip TAM to grow to US\textbackslash{}\$91bn by 2030E Exhibit 4: We expect China AI chip self-sufficiency to reach 70\% in 2030E Exhibit 5: Nvidia's 5090 price continues to rise in China}
\end{figure}
\begin{figure}[htbp]
\centering
\includegraphics[width=0.88\linewidth]{figures/fig_100.jpg}
\caption{Figure 1 - Figure 1: Amazon.com, Inc. - Unit Growth vs Shipping Cost Growth (1Q18-2Q26) Reported net sales of \textbackslash{}\$200.6 billion vs guidance range of \textbackslash{}\$194 billion to \textbackslash{}\$199 billion (+16\%-19\% YOY) came in above the Street \textbackslash{}\$196.8 billion and U}
\end{figure}
{\small\color{refgray}本节综合 9 篇研究；涉及 Bernstein、JPM、MS、NOM、UBS。}
\newpage
\subsection{AI驱动半导体设备与存储：供需重构与盈利上行}
\textbf{AI需求正从云端算力向存储、载板与设备端传导，存储行业由价格周期转向订单周期，设备与载板盈利上行早于预期，供需紧张或延续至2028年。}
\subsubsection*{主流共识}
\begin{itemize}
  \item AI需求持续拉动存储、先进封装与设备市场，2026-2027年供需缺口扩大。
  \item 长期协议（LTA）锁定存储产能，改变行业定价逻辑，供给弹性受限。
  \item HBM4成为存储竞争焦点，三星技术领导力回归，出货量高增。
  \item 半导体设备毛利率上行，东京电子等龙头盈利预测上修。
\end{itemize}
\subsubsection*{分歧与条件差异}
\begin{itemize}
  \item 存储供需紧张持续时间：GS认为2027年缺口大于2026年并延续至2028年，而部分市场观点担忧高利润刺激供给快速释放。
  \item ABF载板毛利率目标：JPM认为60\%经常性水平可在2027下半年实现，但市场此前视为远期目标，时间点分歧明显。
  \item 射频半导体：Citi对Skyworks与Qorvo盈利预测分化，反映合并交易下短期融资成本与长期协同效应的不同权重。
\end{itemize}
\subsubsection*{机构观点}
\paragraph{JPM} 欣兴电子ABF载板毛利率中期目标60\%，预计2027下半年实现，较市场预期提前。驱动因素从涨价转向AI产品组合升级，AI占ABF营收比重从上半年60\%提升至下半年70\%以上，新AI加速器ASP为消费级2-3倍。EMIB-T封装基板价值量约为CoWoS方案2倍，公司有望获关键设备1/3份额。
{\small\color{refgray}数据：2Q26毛利率24.8\%，环比+684bps，同比+1173bps\par}
{\small\color{refgray}数据：3Q26营收环比+15\%，毛利率再扩张6-7个百分点\par}
{\small\color{refgray}数据：AI占ABF营收：上半年60\%→下半年70\%+\par}
{\small\color{refgray}数据：新AI加速器ASP为消费级2-3倍\par}
{\small\color{refgray}边际变化：管理层首次给出ABF毛利率长期锚点60\%，且时间点早于预期；驱动因素从涨价切换至产品组合升级。\par}
\paragraph{MS} 三星电子2Q26业绩超预期，存储业务叙事从周期反弹转向技术领导力回归。HBM4率先推向市场，HBM4E已送样，3Q26 HBM出货量环比增3倍，明年HBM4占HBM销售约60\%。长期协议可能锁定60-70\%存储产能，改变2027年供给弹性。
{\small\color{refgray}数据：2Q26营收172万亿韩元，同比+130\%，环比+28\%\par}
{\small\color{refgray}数据：半导体收入同比+357\%，存储+471\%\par}
{\small\color{refgray}数据：LTA覆盖存储产能60-70\%\par}
{\small\color{refgray}数据：2026e P/E 3.9倍，P/B 2.1倍\par}
{\small\color{refgray}边际变化：存储行业从现货价格主导转向长约订单驱动，2027年供给紧张可能比现货价格显示的更持久。\par}
\paragraph{Bernstein} 东京电子1Q毛利率46.8\%超预期，盈利上行早于市场预期。管理层确认提价自4Q起体现，下个财年初毛利率可达50\%以上，较市场预期提前近一年。2027年WFE市场预期从1700亿上调至1900亿美元，AI之外传统DRAM和CPU也贡献增长。
{\small\color{refgray}数据：1Q收入7324亿日元，同比+33\%\par}
{\small\color{refgray}数据：毛利率46.8\%，营业利润率28.9\%\par}
{\small\color{refgray}数据：FY27/3营收预测上调6\%，FY28/3上调20\%\par}
{\small\color{refgray}数据：目标价从59,200上调至79,300日元，隐含52\%上行空间\par}
{\small\color{refgray}边际变化：毛利率50\%以上时间点提前近一年，盈利预测上修主要来自利润率假设修正而非估值扩张。\par}
\paragraph{Citi} Skyworks与Qorvo近六季度业绩均超预期，但估值倍数被下调，反映射频行业在智能手机平淡下的重新定价。合并交易时间表提前至CY26年内，最早9月底完成。Qorvo FY27毛利率可超50\%，EPS上调10\%；Skyworks因融资利息支出，CY26/27 EPS下调2-3\%。
{\small\color{refgray}数据：Qorvo FY27毛利率>50\%，EPS>7美元\par}
{\small\color{refgray}数据：Skyworks CY26/27 EPS下调2\%/3\%\par}
{\small\color{refgray}数据：SWKS目标价77→64美元，QRVO 100→95美元\par}
{\small\color{refgray}数据：估值倍数从14x/13x降至12x\par}
{\small\color{refgray}边际变化：合并交易时间表提前，短期融资成本侵蚀Skyworks利润，而Qorvo盈利改善建立在交易顺利推进假设上。\par}
\paragraph{GS} 三星电子2Q26业绩符合预告，基本面未因股价回调改变。存储需求持续超供给，DRAM营业利润率接近80\%，NAND超60\%。2027年供需缺口预计比2026年更大，紧张或延续至2028年。长期协议覆盖计划产能60-70\%，HBM收入环比增超80\%，全年HBM销售增长3倍以上。
{\small\color{refgray}数据：2Q26营收171.5万亿韩元，营业利润89.5万亿韩元\par}
{\small\color{refgray}数据：DRAM营业利润率\textasciitilde{}80\%，NAND>60\%\par}
{\small\color{refgray}数据：库存仅2-3周，低于正常水平\par}
{\small\color{refgray}数据：HBM收入环比+80\%+，全年销售增长3倍+\par}
{\small\color{refgray}边际变化：2027年供需缺口比2026年更大，LTA锁定六至七成产能，供给弹性受限，存储价格支撑时间或超预期。\par}
\paragraph{GS} 东京电子1Q营业利润2114亿日元符合预期，但毛利率46.8\%超预期。上半年营业利润指引从4310亿上调至4580亿日元，全年营业利润1万亿日元进入视野，高于彭博一致预期不到9500亿日元。WFE市场2026年增长20\%+至1500亿美元，2027年超1900亿美元。先进封装与DRAM布线设备预期上调。
{\small\color{refgray}数据：1Q营业利润2114亿日元，毛利率46.8\%\par}
{\small\color{refgray}数据：FY3/27营业利润预测上调至10224亿日元\par}
{\small\color{refgray}数据：WFE市场2026年>1500亿美元，2027年>1900亿美元\par}
{\small\color{refgray}数据：存货天数从109天降至84天\par}
{\small\color{refgray}边际变化：全年营业利润预测首次突破1万亿日元，较共识高近8\%；先进封装和DRAM布线设备预期较三个月前上修。\par}
\subsubsection*{关键数据}
\begin{itemize}
  \item ABF载板毛利率目标60\%，2027下半年实现（JPM）
  \item 存储LTA覆盖产能60-70\%，2027年供给弹性受限（MS/GS）
  \item 三星2Q26存储收入同比+471\%，HBM出货3Q环比增3倍（MS）
  \item 东京电子毛利率50\%以上时间点提前至FY3/28初期（Bernstein/GS）
  \item WFE市场2027年预期上调至1900亿美元（Bernstein/GS）
  \item 三星DRAM营业利润率\textasciitilde{}80\%，NAND>60\%（GS）
  \item Qorvo FY27毛利率>50\%，EPS上调10\%（Citi）
  \item Skyworks因合并融资利息支出，CY26 EPS下调2-3\%（Citi）
\end{itemize}
\begin{figure}[htbp]
\centering
\includegraphics[width=0.88\linewidth]{figures/fig_079.jpg}
\caption{Figure 1 - Figure 1: Unimicron ABF GM trend \%}
\end{figure}
\begin{figure}[htbp]
\centering
\includegraphics[width=0.88\linewidth]{figures/fig_108.jpg}
\caption{EXHIBIT 4 - EXHIBIT 5: Tokyo Electron: historical 1-year forward P/E multiple EXHIBIT 6: TEL Q1 revenue was ¥732.4bn, +33\% yoy.}
\end{figure}
\begin{figure}[htbp]
\centering
\includegraphics[width=0.88\linewidth]{figures/fig_178.jpg}
\caption{Exhibit 1 - Exhibit 1: SEC OP waterfall chart (W tn) Exhibit 2: SEC quarterly OP Exhibit 3: SEC HBM revenue and as \% of total DRAM}
\end{figure}
\begin{figure}[htbp]
\centering
\includegraphics[width=0.88\linewidth]{figures/fig_181.jpg}
\caption{Exhibit 3 - Exhibit 6: SEC DRAM key assumptions}
\end{figure}
{\small\color{refgray}本节综合 6 篇研究；涉及 JPM、MS、Bernstein、Citi、GS。}
\newpage
\subsection{AI算力驱动下的被动元件与CCL：供需收紧与高端化共振}
\textbf{AI服务器需求正同时收紧MLCC与高端CCL的供需格局，订单能见度拉长至2027年，产品结构向M7/M8/M9及高容值MLCC迁移，量价齐升与长约签订成为行业新常态。}
\subsubsection*{主流共识}
\begin{itemize}
  \item AI服务器需求是MLCC与CCL供需收紧的核心驱动，订单能见度已延伸至2027年。
  \item 高端产品（M7以上CCL、高容值MLCC）占比提升，推动行业毛利率结构性上行。
  \item 原材料成本上升（玻纤、铜箔）与供应偏紧，支撑CCL与MLCC价格进一步上调。
  \item 客户开始签订长期供货协议（LTA）以锁定供应，行业盈利可见度提升。
\end{itemize}
\subsubsection*{分歧与条件差异}
\begin{itemize}
  \item CCL供应紧张节奏：Citi认为从3Q26开始并延续至2027年，而GS强调台燿提前扩产可能缓解部分缺口，但EMC新产能要到2026年底才上线，供给时间差存在分歧。
  \item MLCC定价策略：三星电机明确表示将更积极管理价格，而日本厂商（村田、TDK等）仍保持观望，亚洲厂商与日系厂商的定价行为存在差异。
  \item 高端CCL毛利率持续性：GS认为M9 CCL毛利率可达50-55\%，但Citi指出TUC高端占比爬坡较慢（39\% vs EMC 60-65\%），高端化斜率不同影响盈利兑现节奏。
  \item 估值倍数调整：Citi将国巨目标价下调（因行业估值倍数从40x降至35x），而MS维持较高目标价（基于2028年31倍PE），对估值中枢判断不一。
\end{itemize}
\subsubsection*{机构观点}
\paragraph{MS} 国巨2Q26业绩超预期，订单出货比从1倍出头跳升至2.2倍，客户开始洽谈长期供货协议，需求能见度显著拉长。3Q26产能利用率预计超过90\%，标准品与高端品同步提升，毛利率有望继续扩张。AI相关需求正在收紧MLCC供给，高容值产品供应趋紧，渠道库存持续消化。公司正从被动元件向主动元件转型，长期方向值得关注。
{\small\color{refgray}数据：订单出货比从1倍出头升至2.2倍\par}
{\small\color{refgray}数据：2Q26营收445亿台币，同比+36\%\par}
{\small\color{refgray}数据：毛利率38.5\%，同比+290bps\par}
{\small\color{refgray}数据：净利率21.2\%，同比+5.9ppt\par}
{\small\color{refgray}边际变化：订单出货比三个月内翻倍，客户行为从短期采购转向长约锁定，需求判断较三个月前明显更积极。\par}
\paragraph{Citi} CCL供应紧张将从3Q26开始显现并延续至2027年，台光电子（EMC）与联茂电子（TUC）2Q26业绩均超预期，管理层不排除3Q26继续调价。EMC高端产品（M7以上）占比达60-65\%，毛利率33.9\%超预期，3Q26指引营收环比+15\%、毛利率35\%±1.5ppt。TUC高端占比39\%，毛利率29.7\%，高端化斜率较缓。行业产能利用率高位，价格调整具备条件。
{\small\color{refgray}数据：EMC 2Q26毛利率33.9\%，环比+4.4ppt\par}
{\small\color{refgray}数据：EMC M7以上产品占比60-65\%\par}
{\small\color{refgray}数据：EMC 3Q26营收指引环比+15\%\par}
{\small\color{refgray}数据：TUC 2Q26毛利率29.7\%，环比+4.6ppt\par}
{\small\color{refgray}边际变化：CCL供应紧张预期从3Q26开始，较此前更明确，且管理层不排除继续调价，价格弹性尚未完全计入指引。\par}
\paragraph{Citi} 国巨2Q26净利润94.18亿台币，同比+88\%，超市场预期3\%。AI收入占比升至16\%，订单出货比2.2，3Q26产能利用率预计突破90\%，供需收紧信号明确。毛利率38.5\%略低于共识，因大宗品占比提升稀释组合，但营业利润率26.4\%超预期，规模效应释放。上调2026-28年EPS预测11-13\%，ASP年复合增速假设上修至22-27\%，目标价因行业估值倍数下调至NT\$1,450。
{\small\color{refgray}数据：2Q26净利润NT\$94.18亿，同比+88\%，环比+18\%\par}
{\small\color{refgray}数据：AI收入占比16\%（1Q26为14-15\%）\par}
{\small\color{refgray}数据：BB ratio从1.3升至2.2\par}
{\small\color{refgray}数据：毛利率38.5\%，同比+2.9ppt\par}
{\small\color{refgray}边际变化：AI收入占比持续提升，BB ratio跳升，管理层不排除继续调价，盈利上修通道打开，但估值倍数因行业因素下调。\par}
\paragraph{GS} 台光电子（EMC）M9 CCL已在3Q26量产，Nvidia HGX R200 UBB为首个采用M9的板卡，LPU用M9从4Q26出货。载板CCL最快4Q26出货，毛利率40-50\%+。管理层提及2H26可能再次调涨CCL价格，因原材料成本上升和玻纤/铜箔供应偏紧。3Q26指引营收环比+15\%以上，毛利率35\%±1.5ppt，未计入涨价和M9增量。M9+需求2027年预计同比增长两倍以上。
{\small\color{refgray}数据：M9 CCL 3Q26量产，1H26已贡献低个位数营收占比\par}
{\small\color{refgray}数据：Nvidia HGX R200 UBB为首个采用M9的板卡\par}
{\small\color{refgray}数据：载板CCL毛利率40-50\%+（公司整体34\%）\par}
{\small\color{refgray}数据：3Q26营收指引环比+15\%以上\par}
{\small\color{refgray}边际变化：M9量产启动标志高端CCL进入新阶段，载板CCL打开第二增长曲线，涨价预期强化盈利弹性。\par}
\paragraph{GS} 日本6月CCL出口量同比+15\%，出口额同比+46\%，量价差约30个百分点，反映产品结构向高端迁移。4-6月季度出口额327亿日元，同比+41\%，环比+12\%。AI服务器封装基板面积更大、层数更多，加上LPDDR和e-SSD需求强劲，拉动CCL价值量上行。三菱瓦斯化学、Resonac、三井金属和日东纺绩均受益，泰国新产能投产（三菱瓦斯化学2025年12月，产能+50\%）确认全球需求韧性。
{\small\color{refgray}数据：6月CCL出口量同比+15\%，出口额同比+46\%\par}
{\small\color{refgray}数据：4-6月出口额327亿日元，同比+41\%，环比+12\%\par}
{\small\color{refgray}数据：三菱瓦斯化学泰国新厂产能+50\%\par}
{\small\color{refgray}数据：三菱瓦斯化学全年BT材料销售目标+15\%\par}
{\small\color{refgray}边际变化：量价差持续扩大，高端化趋势明确，泰国新产能投产增加全球供应，但需求韧性仍待验证。\par}
\paragraph{GS} 日本6月MLCC出口均价环比+4\%，出口量环比+6\%，出口额环比+11\%，同比分别+23\%、+6\%、+31\%，量价齐升显示供需紧张加深。AI和数据中心需求强劲，其他应用领域供给担忧上升，客户正为2026年备货并增加2027年长约请求，需求能见度延伸至两年以上。即将公布的日系厂商财报中，关注BB ratio、定价策略和产能利用率变化。
{\small\color{refgray}数据：6月MLCC出口均价环比+4\%，出口量环比+6\%，出口额环比+11\%\par}
{\small\color{refgray}数据：同比：均价+23\%，量+6\%，额+31\%\par}
{\small\color{refgray}数据：客户开始锁定2026年产能并洽谈2027年长约\par}
{\small\color{refgray}边际变化：量价齐升组合罕见，客户行为从短期补库转向中长期规划，供需紧张程度较5月提升。\par}
\paragraph{GS} 三星电机2Q26营业利润4400亿韩元，超市场预期，MLCC收入环比+19\%、同比+44\%，AI数据中心和汽车电子为主要驱动。管理层预计AI服务器MLCC需求继续增长，1005 47μF产品明年需求可能增加7倍，供应端更紧张，公司将更积极管理定价。FC-BGA下半年产能利用率接近满负荷，已与多数AI/数据中心客户洽谈长期协议，MLCC已签10个长约，盈利可见度提升。
{\small\color{refgray}数据：2Q26营业利润4400亿韩元，超预期\par}
{\small\color{refgray}数据：MLCC收入环比+19\%，同比+44\%\par}
{\small\color{refgray}数据：1005 47μF产品明年需求预计+7倍\par}
{\small\color{refgray}数据：MLCC已签10个长期协议\par}
{\small\color{refgray}边际变化：公司首次明确表示将更主动管理定价，MLCC长约签订数量增加，盈利可见度显著提升。\par}
\paragraph{GS} 台燿科技泰国厂新产能爬坡提前两个月至7月中旬启动，2028年泰国和大陆新增195万张产能（约为现有产能80\%）。M7/M8营收占比从36\%升至39\%，高端化加快。公司正与多家客户推进M8以上等级AI新项目认证，覆盖GPU/AI CPU/ASIC。竞争对手EMC要到2026年底才有新产能，台燿是产能空窗期的主要受益者。800G/1.6T交换机CCL需求2027年弹性高，LEO卫星升级带来增量。
{\small\color{refgray}数据：泰国厂爬坡提前两个月（7月中旬 vs 原计划9月）\par}
{\small\color{refgray}数据：2028年新增产能195万张，约为现有产能80\%\par}
{\small\color{refgray}数据：M7/M8营收占比39\%（1Q26为36\%）\par}
{\small\color{refgray}数据：EMC新产能2026年底才上线\par}
{\small\color{refgray}边际变化：扩产节奏超预期，高端客户认证推进，供给时间差带来份额提升窗口。\par}
\subsubsection*{关键数据}
\begin{itemize}
  \item 国巨订单出货比从1倍出头升至2.2倍
  \item EMC M7以上产品占比60-65\%，毛利率33.9\%
  \item M9 CCL量产启动，毛利率50-55\%+
  \item 日本6月MLCC出口均价环比+4\%，量+6\%
  \item 三星电机MLCC收入环比+19\%，同比+44\%
  \item 台燿泰国厂爬坡提前两个月，M7/M8占比39\%
  \item 日本4-6月CCL出口额同比+41\%
  \item 国巨AI收入占比升至16\%
\end{itemize}
\begin{figure}[htbp]
\centering
\includegraphics[width=0.88\linewidth]{figures/fig_129.jpg}
\caption{Figure 1 - Figure 1. Yageo - 2Q26 earnings review © 2026 Citi Inc. No redistribution without Citi's written permission. Figure 2. Yageo - 2Q26 Sales mix by product © 2026 Citi Inc. No redistribution without Citi's written permission. Figure 3. Yageo - 2Q26 Sales mix by a}
\end{figure}
\begin{figure}[htbp]
\centering
\includegraphics[width=0.88\linewidth]{figures/fig_125.jpg}
\caption{Figure 3 - © 2026 Citi Inc. No redistribution without Citi's written permission. © 2026 Citi Inc. No redistribution without Citi's written permission. Figure 3. EMC – 12M Forward P/E Band © 2026 Citi Inc. No redistribution without Citi's written permission. Figure 4. EMC}
\end{figure}
\begin{figure}[htbp]
\centering
\includegraphics[width=0.88\linewidth]{figures/fig_177.jpg}
\caption{Exhibit 1 - Exhibit 1: MLCC export volume and average export price (ASP) \#\# Price target risks and methodology \#\# Price Target Risks and Methodology - Murata Mfg. Valuation methodology: We are Buy-rated on Murata Mfg. with a 12-month price t}
\end{figure}
\begin{figure}[htbp]
\centering
\includegraphics[width=0.88\linewidth]{figures/fig_183.jpg}
\caption{Exhibit 1 - Exhibit 1: SEMCO shares are trading at 30.6X 12m FWD P/E SEMCO 12m FWD P/E Exhibit 2: SEMCO shares are trading at 3.6X 12m FWD P/B with 2026E/2027E ROEs of 13\%/18\% SEMCO 12m FWD P/B vs. ROE Exhibit 3: SEMCO earnings revisions}
\end{figure}
{\small\color{refgray}本节综合 8 篇研究；涉及 MS、Citi、GS。}
\newpage
\subsection{消费与零售：高端韧性延续，大众消费分化加剧}
\textbf{全球消费与零售板块呈现明显分化：高端品牌（爱马仕、法拉利）凭借定价权与个性化配置维持盈利韧性，大众消费（啤酒、运动服饰）则在中国市场承压，但新业态与产品结构升级提供局部亮点。}
\subsubsection*{主流共识}
\begin{itemize}
  \item 高端品牌定价权稳固，个性化配置与产品组合升级持续驱动单价提升，销量增长放缓不碍利润率扩张。
  \item 中国市场消费信心恢复缓慢，大众消费品类（啤酒、运动服饰）销量承压，但产品结构升级（高端化、O2O渠道）仍在推进。
  \item 供应链扰动（如罗技供应商事件）是短期增长的主要制约，需求端趋势整体健康。
\end{itemize}
\subsubsection*{分歧与条件差异}
\begin{itemize}
  \item 爱马仕增长可持续性：BARC认为需看到增速稳定回到高个位数区间的明确证据，而市场对产能扩张与需求匹配度存在分歧。
  \item 百威亚太中国区复苏时点：Bernstein预计下半年利润率继续承压，GS则认为均价正增长与O2O渠道高增显示高端化逻辑未破，底部或已临近。
  \item adidas鞋履与服装增速差（35\% vs 1\%）的持续性：GS认为鞋履订单簿改善将带动4Q26/1Q27回升，但行业性折扣不确定性仍存。
\end{itemize}
\subsubsection*{机构观点}
\paragraph{BARC} 爱马仕2Q26收入41亿欧元，恒定汇率增长6.7\%符合预期，但皮具部门增长10.2\%低于11-12\%长期目标区间，亚太区（除日本）仅增2.5\%成为主要拖累。上半年EBIT利润率41\%超预期，推动EPS上调约1\%，但增长可持续性仍是核心疑虑。产能扩张计划对应6-7\%销量增长，供给端到2030年不构成约束，但Q3预期增速7.6\%仍不足以证明增长回到高个位数。估值吸引力渐增，但需等待更明确的增长证据。
{\small\color{refgray}数据：2Q26收入41亿欧元，cc增长6.7\%\par}
{\small\color{refgray}数据：皮具部门增长10.2\%，低于预期10.7\%\par}
{\small\color{refgray}数据：亚太区（除日本）增长2.5\%，低于预期3.2\%\par}
{\small\color{refgray}数据：香水美妆负增长9.5\%，远低于预期-1.5\%\par}
{\small\color{refgray}边际变化：全年增长预期下调，但利润率预期上调，反映增长与盈利的再平衡。\par}
\paragraph{Bernstein} 百威亚太2Q26中国区销量同比下滑9.7\%，较Q1的-1.5\%显著扩大，价格组合转正1.2\%提供部分缓冲，但cc收入仍降8.6\%。韩国区在低基数下实现低双位数销量增长，EBITDA同比增长26.3\%，利润率扩张413个基点。下调2026年亚太西部EBITDA增速预期至-8.0\%，上调东部至+5.7\%，整体cc EBITDA预期下调1.9\%。中国区仍处投入期，渠道扩张与运营去杠杆使利润率承压，预计下半年亚太西部利润率再调整100个基点。
{\small\color{refgray}数据：中国区销量同比-9.7\%（Q1为-1.5\%）\par}
{\small\color{refgray}数据：价格组合+1.2\%\par}
{\small\color{refgray}数据：韩国区EBITDA同比+26.3\%\par}
{\small\color{refgray}数据：2026年亚太西部EBITDA增速预期下调至-8.0\%\par}
{\small\color{refgray}边际变化：中国区销量降幅加速扩大，但价格组合转正显示高端化仍在推进。\par}
\paragraph{GS} 百胜中国2Q26同店销售增长改善，管理层对3Q26给出正增长指引，新业态成为第二增长曲线：K Coffee Cafe年底达5000家，K Pro开店指引上调至800家，必胜客汉堡站超200家且对母店有双位数销售提升。Pizza Hut收购预计2026年8月完成，将带来利润率结构变化。上调2026-28年盈利预测约2.5\%，2027-28年开店指引从每年600家上调至800家以上。外卖平台补贴战趋于理性，行业连锁化率提升，头部品牌受益。
{\small\color{refgray}数据：K Coffee Cafe年底5000家，销售额弹性至20亿元\par}
{\small\color{refgray}数据：K Pro开店指引上调至800家\par}
{\small\color{refgray}数据：必胜客汉堡站超200家，对母店双位数销售提升\par}
{\small\color{refgray}数据：2026-28年盈利预测上调约2.5\%\par}
{\small\color{refgray}边际变化：新业态从试验性布局转向可量化的收入来源，开店指引大幅上调。\par}
\paragraph{JPM} 罗技国际F1Q27营收12.27亿美元，同比增长7\%，超预期约2\%；剔除6100万美元关税退税后毛利率44.8\%，仍高于预期。增长由鼠标（+16\%）、游戏（+12\%）和视频协作（+11\%）拉动，北美回到中个位数增长。但6月末半导体供应商晶圆厂临时关闭，影响游戏和个人工作空间产品线，管理层估计F2Q影响2000万美元、F3Q影响2亿美元，F4Q基本消化。需求端趋势健康，问题在供给端，市场应关注供应恢复时点而非需求走弱。
{\small\color{refgray}数据：F1Q27营收12.27亿美元，同比+7\%\par}
{\small\color{refgray}数据：毛利率44.8\%（剔除关税退税后）\par}
{\small\color{refgray}数据：鼠标+16\%，游戏+12\%，视频协作+11\%\par}
{\small\color{refgray}数据：供应商事件影响：F2Q 2000万美元，F3Q 2亿美元\par}
{\small\color{refgray}边际变化：业绩超预期但供应链扰动成为焦点，供给冲击有时间表而非需求转弱。\par}
\paragraph{GS} 百威亚太2Q26集团均价上涨2\%与销量下滑4.1\%并行，显示产品结构升级仍在推进，但量的不确定性未解除。中国区O2O渠道高双位数增长，由高端SKU和包装创新驱动，但促销投入较高。管理层将稳定中国销量和恢复市场份额列为2026下半年首要任务，未给出V型反转时间表。下调2026年中国区有机销售至-3.5\%（销量-5\%，均价+1.5\%），Q3预期销量同比-8\%，Q4在低基数上预计销量-1\%、均价+5\%、EBITDA+8\%。原材料成本压力（铝价同比+19\%）将在下半年至2027年显现。
{\small\color{refgray}数据：集团均价+2\%，销量-4.1\%\par}
{\small\color{refgray}数据：中国区O2O渠道高双位数增长\par}
{\small\color{refgray}数据：2026年中国区有机销售预期下调至-3.5\%\par}
{\small\color{refgray}数据：Q3预期销量-8\%，Q4预期销量-1\%、均价+5\%\par}
{\small\color{refgray}边际变化：均价正增长与销量负增长并行，显示消费者未降级但饮用量减少，高端化逻辑仍在。\par}
\paragraph{GS} adidas 2Q26全球cFX销售增长14\%，大中华区增长15\%，但品类分化显著：服装加速至35\%，鞋履仅增1\%，增速差达34个百分点。对供应链含义不均：服装代工厂（申洲国际）订单动能强于鞋履代工厂（华利、裕元），分销商（滔搏、宝胜）报表改善但传导有时滞，本土品牌（安踏、李宁、特步）面临竞争压力。管理层预计鞋履增长在4Q26/1Q27改善，因批发订单簿好转和渠道库存正常化，但生活方式鞋类折扣不确定性未消。库存增加12\%，结构健康，过季产品占比仅9\%。
{\small\color{refgray}数据：adidas全球cFX销售+14\%，大中华区+15\%\par}
{\small\color{refgray}数据：服装+35\%，鞋履+1\%\par}
{\small\color{refgray}数据：Nike中国区-17\%，宝胜-3\%\par}
{\small\color{refgray}数据：库存+12\%，过季产品占比9\%\par}
{\small\color{refgray}边际变化：服装与鞋履增速差扩大至34个百分点，供应链受益分层明显。\par}
\paragraph{GS} 法拉利2Q26交付3,366辆，低于一致预期1.9\%，但收入19.38亿欧元超预期4.3\%，调整后EBIT超预期6.3\%。核心驱动力是个性化配置率升至20\%，高于中期指引18\%，价格组合贡献1.22亿欧元（GS估计0.82亿）。订单簿已覆盖2027全年，上调2026年指引：收入至76亿欧元，调整后EBIT至不低于22.6亿欧元。当前指引隐含26-30年收入CAGR超5\%，市场开始将5\%视为增长基础而非上限。下半年F80超跑和296 Versione Speciale放量，定制比例仍有提升空间。
{\small\color{refgray}数据：交付3,366辆，低于预期1.9\%\par}
{\small\color{refgray}数据：收入19.38亿欧元，超预期4.3\%\par}
{\small\color{refgray}数据：个性化配置率20\%，高于中期指引18\%\par}
{\small\color{refgray}数据：价格组合贡献1.22亿欧元（GS估计0.82亿）\par}
{\small\color{refgray}边际变化：盈利逻辑从销量驱动转向单价驱动，个性化配置成为定价权核心指标。\par}
\subsubsection*{关键数据}
\begin{itemize}
  \item 爱马仕2Q26收入41亿欧元，cc增长6.7\%，亚太区（除日本）仅增2.5\%
  \item 百威亚太中国区销量同比-9.7\%，价格组合+1.2\%
  \item 百胜中国K Coffee Cafe年底5000家，K Pro开店指引上调至800家
  \item 罗技F1Q27营收12.27亿美元，供应商事件影响F3Q 2亿美元
  \item adidas服装+35\% vs 鞋履+1\%，增速差34个百分点
  \item 法拉利个性化配置率20\%，价格组合贡献1.22亿欧元
  \item 法拉利订单簿覆盖2027全年，2026年收入指引上调至76亿欧元
\end{itemize}
\begin{figure}[htbp]
\centering
\includegraphics[width=0.88\linewidth]{figures/fig_001.jpg}
\caption{FIGURE 1 - Upside/Downside scenarios EQUAL WEIGHT \#\# Why EQUAL WEIGHT? Assume single-digit top-line growth as other segments grow at mid-single digits, diluting the higher structural growth of the Leather Goods division. EBIT margin remains in the 36-37\% range. Terminal }
\end{figure}
\begin{figure}[htbp]
\centering
\includegraphics[width=0.88\linewidth]{figures/fig_031.jpg}
\caption{EXHIBIT 4 - EXHIBIT 6: We have increased East cc EBITDA growth to \$+5.7\textbackslash{}\%\$ but cut West's by \$3.5\textbackslash{}\%\$ and now we expect overall cc EBITDA growth to decline \$4.7\textbackslash{}\%\$ in FY26 Bernstein Estimates: Bud APAC FY26e cc Normalized EBITDA Growth YoY\%}
\end{figure}
\begin{figure}[htbp]
\centering
\includegraphics[width=0.88\linewidth]{figures/fig_048.jpg}
\caption{Exhibit 2 - Exhibit 3: We expect KFC SSSG/cFX restaurant profits to be +1\%/+4\% yoy respectively in 3Q26E... KFC quarterly SSSG vs. quarterly restaurants profit growth Exhibit 4: ...and +1\%/+5\% yoy SSSG/cFX restaurant profits change for Pizza}
\end{figure}
\begin{figure}[htbp]
\centering
\includegraphics[width=0.88\linewidth]{figures/fig_170.jpg}
\caption{Exhibit 1 - Exhibit 1: \#\# 2Q26 results vs. GSe and consensus Exhibit 2: Ferrari 2Q26 results vs. GSe and Visible Alpha consensus}
\end{figure}
{\small\color{refgray}本节综合 9 篇研究；涉及 BARC、Bernstein、GS、JPM。}
\newpage
\subsection{医疗健康：临床测序需求回暖，中国创新药进入全球后期开发}
\textbf{医疗健康板块呈现结构性分化：临床测序需求从替换转向增量，CRO受益于新兴生物制药融资回暖，而中国创新药资产正获得跨国药企大规模后期临床投入，全球商业化价值加速兑现。}
\subsubsection*{主流共识}
\begin{itemize}
  \item 临床测序需求正从平台替换转向增量扩张，NovaSeq X装机中约70\%来自临床客户，其中一半为新增需求。
  \item 新兴生物制药融资环境显著改善，2Q26融资规模达350亿美元，为去年同期的两倍以上，直接推动CRO订单增长。
  \item 跨国药企对中国创新药资产的投入已从许可引进转向大规模后期临床开发，GSK和AZ正为引进的ADC和抗体项目规划多项III期研究。
  \item CRO行业普遍上调全年收入指引，反映生物制药研发外包需求回暖，EBP客户成为主要驱动力。
\end{itemize}
\subsubsection*{分歧与条件差异}
\begin{itemize}
  \item 临床测序耗材收入是否会出现断崖：Bernstein认为切换阵痛期已过，耗材增长15\%证明量价平衡；但部分市场观点仍担忧单样本价格下降的长期影响。
  \item AI对CRO行业的影响方向存在分歧：管理层认为AI将在2027年后增加CRO需求，但研报提示AI对商业咨询板块可能构成压力，临床前与商业分析受益程度不同。
  \item CRO公司信用质量分化：IQVIA和Avantor上调指引后信用面改善，但Fortrea高杠杆（总杠杆6.7倍）仍是风险点，即便EBITDA利润率提升。
  \item 中国创新药资产价值兑现节奏：GSK和AZ的III期规划显示信心，但BLA提交和商业化时间点（如2028年）存在执行不确定性。
\end{itemize}
\subsubsection*{机构观点}
\paragraph{Bernstein} Illumina 2Q26业绩超预期，营收11.59亿美元（高于预期4\%），EPS 1.31美元（高于预期7\%），全年指引上调至46.0-46.4亿美元。NovaSeq X单季装机超95台，环比一季度83台显著提升，其中70\%来自临床客户，且临床装机中约50\%为增量需求而非替换。临床耗材收入增长15\%（美国加拿大超20\%），X平台已占高通量耗材收入59\%，此前担忧的'临床耗材断崖'正被逐季数据削弱。装机对收入贡献主要在FY27，因临床客户需6-9个月达到正常耗材消耗。
{\small\color{refgray}数据：营收11.59亿美元，超预期4\%\par}
{\small\color{refgray}数据：EPS 1.31美元，超预期7\%\par}
{\small\color{refgray}数据：NovaSeq X单季装机95台，环比+14\%\par}
{\small\color{refgray}数据：临床装机占比70\%，其中50\%为增量\par}
{\small\color{refgray}边际变化：市场此前担忧临床客户从NovaSeq 6000向X迁移导致耗材收入断崖，但2Q26数据显示增量需求占比高，切换阵痛期已过，耗材增长曲线持续兑现。\par}
\paragraph{JPM} GSK和AZ在2Q26业绩中披露对中国创新药资产的后期开发计划：GSK为翰森制药的B7-H3 ADC（Ris-Rez）和B7-H4 ADC（Mo-Rez）规划多项III期研究，预计到2028年分别运行6项和5项III期，并计划2028年提交Ris-Rez二线小细胞肺癌BLA；GSK还启动恒瑞长效抗TSLP单抗（GSK'283）在哮喘、慢阻肺、鼻息肉三个适应症的III期，恒瑞可获得超10亿美元里程碑付款。AZ的CLDN18.2 ADC（Sone-Ve）在全球III期CLARITY-Gastric01达到总生存期终点，成为首个显示OS获益的CLDN18.2 ADC。这些进展表明中国资产已进入跨国药企全球核心管线。
{\small\color{refgray}数据：GSK计划2028年运行6项Ris-Rez III期和5项Mo-Rez III期\par}
{\small\color{refgray}数据：Ris-Rez预计2028年提交2L小细胞肺癌BLA\par}
{\small\color{refgray}数据：GSK'283启动3个适应症III期\par}
{\small\color{refgray}数据：恒瑞里程碑付款超10亿美元\par}
{\small\color{refgray}边际变化：跨国药企对中国创新药资产的投入从许可引进转向大规模后期临床开发，且AZ的CLDN18.2 ADC达到OS终点，验证了中国ADC平台的全球竞争力。\par}
\paragraph{JPM} IQVIA、Fortrea和Avantor三家CRO/制药服务公司2Q26业绩均上调全年收入指引，核心驱动力是新兴生物制药（EBP）融资回暖。IQVIA 2Q26收入44亿美元（同比+8.7\%），全年收入指引中值从5.8\%上调至6.5\%；EBP融资规模达350亿美元，为去年同期两倍以上，EBP占IQVIA研发收入约35\%，且EBP研发支出增速预计为大药企的2-3倍。Fortrea EBITDA利润率从7.3\%升至8.6\%，但总杠杆6.7倍、第一留置权杠杆7.5倍，信用风险仍高。AI对行业影响尚处早期，管理层预计2027年后AI辅助药物发现将增加CRO需求，但商业咨询板块可能承压。
{\small\color{refgray}数据：IQVIA 2Q26收入44亿美元，同比+8.7\%\par}
{\small\color{refgray}数据：IQVIA全年收入指引中值上调至6.5\%\par}
{\small\color{refgray}数据：EBP融资规模350亿美元，同比翻倍\par}
{\small\color{refgray}数据：EBP占IQVIA研发收入35\%\par}
{\small\color{refgray}边际变化：生物制药融资环境改善推动CRO需求回暖，三家CRO同步上调指引；但Fortrea高杠杆与IQVIA稳健增长形成分化，AI影响仍不确定。\par}
\subsubsection*{关键数据}
\begin{itemize}
  \item Illumina 2Q26营收11.59亿美元，超预期4\%；EPS 1.31美元，超预期7\%
  \item NovaSeq X单季装机95台，70\%临床客户，50\%增量需求
  \item 临床耗材收入增长15\%，X平台占高通量耗材收入59\%
  \item IQVIA 2Q26收入44亿美元，EBP融资350亿美元（同比翻倍）
  \item EBP占IQVIA研发收入35\%，研发支出增速为大药企2-3倍
  \item GSK计划2028年运行6项Ris-Rez III期和5项Mo-Rez III期
  \item AZ的CLDN18.2 ADC达到OS终点，首个显示OS获益
  \item Fortrea总杠杆6.7倍，EBITDA利润率升至8.6\%
\end{itemize}
\begin{figure}[htbp]
\centering
\includegraphics[width=0.88\linewidth]{figures/fig_036.jpg}
\caption{EXHIBIT 6 - EXHIBIT 6: Illumina Core Revenues EXHIBIT 7: Illumina Consumables and Instruments Revenues EXHIBIT 8: Illumina Core Operating Margin}
\end{figure}
\begin{figure}[htbp]
\centering
\includegraphics[width=0.88\linewidth]{figures/fig_037.jpg}
\caption{EXHIBIT 6 - EXHIBIT 6: Illumina Core Revenues EXHIBIT 7: Illumina Consumables and Instruments Revenues EXHIBIT 8: Illumina Core Operating Margin EXHIBIT 9: NovaSeq X shipments}
\end{figure}
\begin{figure}[htbp]
\centering
\includegraphics[width=0.88\linewidth]{figures/fig_038.jpg}
\caption{EXHIBIT 7 - EXHIBIT 7: Illumina Consumables and Instruments Revenues EXHIBIT 8: Illumina Core Operating Margin EXHIBIT 9: NovaSeq X shipments \#\# APPENDIX - FINANCIAL FORECASTS}
\end{figure}
\begin{figure}[htbp]
\centering
\includegraphics[width=0.88\linewidth]{figures/fig_039.jpg}
\caption{EXHIBIT 8 - EXHIBIT 8: Illumina Core Operating Margin EXHIBIT 9: NovaSeq X shipments \#\# APPENDIX - FINANCIAL FORECASTS EXHIBIT 10: ILMN Revenue and IS}
\end{figure}
{\small\color{refgray}本节综合 3 篇研究；涉及 Bernstein、JPM。}
\newpage
\subsection{电力设备订单排至2030年代，中国自动化复苏超预期}
\textbf{全球AI数据中心需求正将燃气轮机、气体发动机和SOFC等电力设备订单周期拉长至2030年代，供给紧张持续；中国工业自动化市场在OEM端驱动下加速增长，全年增速预期上调，国产替代与多下游回暖共振。}
\subsubsection*{主流共识}
\begin{itemize}
  \item AI数据中心电力需求强劲，分布式发电与电力设备订单能见度延伸至2030年代，供给紧张持续数年。
  \item 中国工业自动化市场2Q26同比增长6\%，OEM端增长17\%，全年增速预期从+2\%上调至+5\%。
  \item 国产替代加速，工业机器人国产化率突破61\%，伺服、小型PLC等核心产品份额持续提升。
  \item 下游复苏覆盖面扩大，电池、电子、半导体、机器人等行业增长加快，光伏转正。
\end{itemize}
\subsubsection*{分歧与条件差异}
\begin{itemize}
  \item 电力设备需求持续性：JPM认为订单排至2030年代无断崖风险，但部分观点担忧产能扩张后供给过剩。
  \item 中国自动化复苏动力：MS强调AI与设备更新驱动，但项目市场（化工、冶金）仍同比下滑，复苏结构不均衡。
  \item 宏发股份毛利率回升：Citi预计3Q26回升至35\%，依赖银价回落与提价，若银价反弹或提价乏力则存在不确定性。
\end{itemize}
\subsubsection*{机构观点}
\paragraph{JPM} JPM通过交叉比对Bloom Energy、Innio和GE Vernova三家二季度业绩，确认AI数据中心电力需求强劲，分布式发电订单簿饱满，产能排期延伸至2030年代，未见需求断崖或供给过剩。Bloom Energy数据中心订单超FY25全年两倍，客户能见度至2030年；潍柴SOFC爬坡路径明确，2026年30MW、2027年200MW，AIDC发动机2026年出货指引3500-4000台，2027年底产能规划8000-10000台。估值已回调至工业周期股水平，H股和A股分别回调约19\%和13\%，股息率约4\%，大股东增持提供支撑。
{\small\color{refgray}数据：Bloom Energy数据中心订单超FY25全年两倍\par}
{\small\color{refgray}数据：潍柴SOFC 2026年30MW、2027年200MW\par}
{\small\color{refgray}数据：AIDC发动机2026年出货3500-4000台，2027年底产能8000-10000台\par}
{\small\color{refgray}数据：H股回调19\%、A股回调13\%，股息率约4\%\par}
{\small\color{refgray}边际变化：从单一公司订单转向三家不同技术路线交叉验证，订单能见度以年为单位，排除了个别公司波动，指向结构性扩张。\par}
\paragraph{MS} MS上调2026年中国自动化市场增速预期至+5\%（原+2\%），2Q26整体市场同比增长6\%，OEM端增长17\%而项目市场下降1\%。核心FA产品增速普遍加快：伺服+25\%、工业机器人+17\%、小型PLC+15\%、中大PLC+14\%、低压变频+7\%，其中协作机器人增速高达45\%。下游电池、电子、半导体、机器人行业增长加快，光伏转正，但化工、冶金、电梯仍调整。国产化率持续提升：工业机器人61\%（+5ppt）、伺服59\%（+4ppt）、小型PLC 39\%（+4ppt）。复苏由AI及物理AI应用、能源安全、设备更新、装备出海四因素驱动。
{\small\color{refgray}数据：2Q26自动化市场+6\%，OEM+17\%，项目市场-1\%\par}
{\small\color{refgray}数据：2026全年增速预期从+2\%上调至+5\%\par}
{\small\color{refgray}数据：伺服+25\%、机器人+17\%、协作机器人+45\%\par}
{\small\color{refgray}数据：工业机器人国产化率61\%（+5ppt）\par}
{\small\color{refgray}边际变化：复苏从单一行业拉动转为多行业同步回暖，全年增速预期上调，上行周期确认。\par}
\paragraph{Citi} Citi认为宏发股份2Q26收入59亿元（+36\%），超预期4个百分点；毛利率33.2\%，超预期2.3个百分点，但同比回落1.5个百分点。增长呈全品类特征：新能源车高压直流继电器收入36亿元（+45\%）、家电继电器23亿元（+40\%，提价贡献20个百分点）、新能源继电器10亿元（+60\%）、PDU/BDU收入增速超100\%。毛利率有望在3Q26回升至约35\%，驱动来自持续提价与银价回落（银价自高点回落约20\%）。上调2026-2028年净利润预测5\%/11\%/12\%，目标价从43元上调至45元。800V AIDC继电器预计最快4Q26向台达或光宝供货。
{\small\color{refgray}数据：2Q26收入59亿元，+36\%，超预期4个百分点\par}
{\small\color{refgray}数据：毛利率33.2\%，超预期2.3个百分点\par}
{\small\color{refgray}数据：家电继电器提价贡献20个百分点收入增量\par}
{\small\color{refgray}数据：银价自高点回落约20\%\par}
{\small\color{refgray}边际变化：增长逻辑从单一需求故事转向提价与银价成本改善双驱动，毛利率回升预期强化。\par}
\subsubsection*{关键数据}
\begin{itemize}
  \item Bloom Energy数据中心订单超FY25全年两倍，订单能见度至2030年
  \item 潍柴SOFC产能2026年30MW→2027年200MW，AIDC发动机2027年底产能8000-10000台
  \item 中国自动化市场2Q26同比+6\%，OEM+17\%，项目市场-1\%
  \item 2026全年自动化增速预期从+2\%上调至+5\%
  \item 协作机器人增速45\%，工业机器人国产化率61\%（+5ppt）
  \item 宏发2Q26收入59亿元（+36\%），毛利率33.2\%，超预期2.3个百分点
  \item 宏发家电继电器提价贡献20个百分点收入增量，银价回落约20\%
  \item 宏发PDU/BDU收入增速超100\%，800V AIDC继电器最快4Q26供货
\end{itemize}
\begin{figure}[htbp]
\centering
\includegraphics[width=0.88\linewidth]{figures/fig_088.jpg}
\caption{Exhibit 1 - Exhibit 1: MIR forecast – by OEM/project market}
\end{figure}
\begin{figure}[htbp]
\centering
\includegraphics[width=0.88\linewidth]{figures/fig_090.jpg}
\caption{Exhibit 3 - Exhibit 6: China's automation market accelerated to +6\% y-y in 2Q26, mainly driven by OEM market}
\end{figure}
\begin{figure}[htbp]
\centering
\includegraphics[width=0.88\linewidth]{figures/fig_120.jpg}
\caption{Figure 1 - Figure 1. Hongfa: 2Q26/1H26 results snapshot © 2026 Citi Inc. No redistribution without Citi's written permission. Figure 2. Hongfa: Revenue by product and YoY growth, 1H26 © 2026 Citi Inc. No redistribution without Citi's written permission. Figure 3. Hongfa:}
\end{figure}
\begin{figure}[htbp]
\centering
\includegraphics[width=0.88\linewidth]{figures/fig_121.jpg}
\caption{Figure 2 - Figure 2. Hongfa: Revenue by product and YoY growth, 1H26 © 2026 Citi Inc. No redistribution without Citi's written permission. Figure 3. Hongfa: Quarterly GPM vs. silver price Figure 4. Hongfa: Earnings forecast revisions © 2026 Citi Inc. No redistribution wi}
\end{figure}
{\small\color{refgray}本节综合 3 篇研究；涉及 JPM、MS、Citi。}
\newpage
\subsection{能源与公用事业：AI电力叙事重构与高股息防御切换}
\textbf{AI与数据中心驱动的电力需求叙事正从主题转向财报验证，但设备端估值已大幅回撤，资金转向低beta高股息公用事业；电池装车量高增由单车带电量、商用车和高端化结构性驱动，而非销量；美国FCC逆变器限制影响偏中期，现有型号暂安全。}
\subsubsection*{主流共识}
\begin{itemize}
  \item 电网建设与数据中心电力需求进入兑现期，承包商和设备商订单指引上调（如Quanta EBITDA超预期21\%，上调全年指引13\%）。
  \item 高股息公用事业（城燃、水务）在美债收益率上行和AI情绪转弱下具备防御价值，未来30天可能跑赢。
  \item 中国动力电池装车量增速（+33\%）远超销量（+10\%），由单车带电量提升、商用车占比上升和高端化驱动。
  \item 美国FCC逆变器限制对现有型号影响有限，实际冲击推迟至新型号审批，且存在豁免通道。
\end{itemize}
\subsubsection*{分歧与条件差异}
\begin{itemize}
  \item FCC'远程连接'定义宽严未定：若扩大解释，阳光电源新型号逆变器无法在美销售；若按字面解释，影响有限。
  \item 电网设备出口商估值分歧：思源电气共识为增长加速，华明电力装备则是低估值修复，两者预期差方向不同。
  \item 光纤与燃气轮机方向遭质疑：新产能和竞争加剧担忧压制估值，尽管中天科技PE已回落至10倍。
  \item 天然气短期供应宽松（产量指引上调）与长期需求上行（数据中心增量20Bcf/d）并存，价格路径不确定。
\end{itemize}
\subsubsection*{机构观点}
\paragraph{JPM} FCC新规覆盖逆变器和储能PCS，适用于公用事业级和户用离网，比预期严格，但已获授权的现有型号仍可销售，影响推迟。美国市场占阳光电源2025年毛利20-25\%，若完全退出盈利下修20-25\%，但非美储能装机2026-27年CAGR约40\%可对冲。关键变量是'远程连接'定义，公司产品用物理线缆，若FCC扩大解释则新型号受限。当前估值已部分消化不确定性。
{\small\color{refgray}数据：美国市场占阳光电源2025年毛利20-25\%\par}
{\small\color{refgray}数据：若完全退出美国市场，盈利下修20-25\%\par}
{\small\color{refgray}数据：非美储能装机2026-27年CAGR约40\%\par}
{\small\color{refgray}数据：公司报价自6月30日以来跑输大盘约25\%\par}
{\small\color{refgray}边际变化：FCC新规细则落地，范围超预期，但现有型号豁免使影响推迟，市场关注点转向'远程连接'定义和豁免申请。\par}
\paragraph{MS} 上海路演反馈显示，AI与出口驱动的电力设备股自年初高点普遍回撤超50\%，但资金未离场，正重新分配。电网设备出口商共识度最高：思源电气上半年新签订单稳健，预期三季度起增长加速；华明电力装备估值低于20倍2027年PE且股息率约4\%，吸引偏空资金回补。光纤担忧最多（新产能和跨界进入者），中天科技10倍PE但分歧大；燃气轮机因估值较高兴趣减弱。
{\small\color{refgray}数据：电力设备股自2026年初高点回撤超50\%\par}
{\small\color{refgray}数据：思源电气与华明电力装备估值约20倍2027年PE\par}
{\small\color{refgray}数据：华明股息率约4\%\par}
{\small\color{refgray}数据：中天科技PE回落至10倍\par}
{\small\color{refgray}边际变化：资金从光纤、燃气轮机等前期热门方向转向电网设备出口商，估值保护与增长弹性成为筛选核心。\par}
\paragraph{Citi} 上半年中国动力电池装车451.6GWh，同比+33\%，远超销量增速10\%，缺口由单车带电量提升（纯电+12\%、插混+27\%、商用车+14\%）、商用车占比升至23\%（装车量+94\%）和高端车型占比上升（A级从40\%降至30\%）驱动。宁德时代市占率升至45\%，比亚迪降至20\%，两家合计65\%。碳酸锂价格在8-9月旺季前或有支撑。
{\small\color{refgray}数据：上半年动力电池装车451.6GWh，同比+33\%\par}
{\small\color{refgray}数据：新能源汽车销量同比+10\%\par}
{\small\color{refgray}数据：插混单车带电量同比+27\%\par}
{\small\color{refgray}数据：商用车装车占比从16\%升至23\%，装车量+94\%\par}
{\small\color{refgray}边际变化：装车量与销量增速缺口被重新解释为结构性因素（带电量、商用车、高端化），而非库存或统计口径。\par}
\paragraph{Citi} AI相关谨慎情绪和10年期美债收益率升至4.61\%推动资金向低beta公用事业轮动。未来30天高股息标的可能跑赢：ENN Energy股息率6.7\%、CR Gas和北京企业5.9\%、粤海和光大环境5.7\%、新奥天然气5.3\%。过去三个月电网设备跌幅最大（HD现代电气-54\%、晓星重工-50\%、东方电气-47\%），城燃和水务抗跌。城燃公司资本开支收缩、自由现金流转正，有能力维持派息。
{\small\color{refgray}数据：10年期美债收益率从3.96\%升至4.61\%\par}
{\small\color{refgray}数据：ENN Energy股息率6.7\%，CR Gas和北京企业5.9\%\par}
{\small\color{refgray}数据：HD现代电气跌54\%，晓星重工跌50\%，东方电气跌47\%\par}
{\small\color{refgray}数据：新奥能源承诺2026年派息不低于3.0港元\par}
{\small\color{refgray}边际变化：市场从AI电力需求叙事转向高股息防御，设备端估值大幅回撤，运营端盈利兑现尚未被检验。\par}
\paragraph{GS} 二季度财报显示电力需求增长从主题变为可验证的订单和指引。Quanta调整后EBITDA超预期21\%，上调全年有机增长指引13\%，完成四项收购，预计2030年EPS CAGR约19\%。Enphase四季度指引好于预期，固态变压器进入RFI/RFP谈判，对应多GW级数据中心管道。天然气短期供应宽松但长期需求增量约20Bcf/d。炼化利润率超预期（Valero EPS 12.54美元），全球炼厂开工率同比下降6.5mb/d。
{\small\color{refgray}数据：Quanta调整后EBITDA超预期21\%，上调全年指引13\%\par}
{\small\color{refgray}数据：Quanta预计2030年EPS CAGR约19\%\par}
{\small\color{refgray}数据：Enphase固态变压器对应多GW级数据中心管道\par}
{\small\color{refgray}数据：天然气长期需求增量约20Bcf/d\par}
{\small\color{refgray}边际变化：电力需求从叙事进入财报兑现期，利润分配向电力基础设施、天然气管道和本土工业材料倾斜。\par}
\subsubsection*{关键数据}
\begin{itemize}
  \item 中国上半年动力电池装车451.6GWh，同比+33\%，销量仅+10\%
  \item 宁德时代市占率升至45\%，比亚迪降至20\%
  \item 美国FCC限制中国逆变器，阳光电源美国市场占毛利20-25\%
  \item Quanta EBITDA超预期21\%，上调全年指引13\%
  \item 10年期美债收益率升至4.61\%，高股息公用事业股息率5.3-6.7\%
  \item 电力设备股自高点回撤超50\%，思源电气和华明PE约20倍
  \item 全球炼厂开工率同比下降6.5mb/d
\end{itemize}
\begin{figure}[htbp]
\centering
\includegraphics[width=0.88\linewidth]{figures/fig_116.jpg}
\caption{Figure 1 - \#\# CITI'S TAKE What's going on YTD - According to Real Lithium, vehicle battery installation in China totaled 451.6 GWh, +33\% YoY. CATL's (excl CATL-SAIC JV) market share was up by 3ppt YoY to 45\% during 6M26, whilst BYD's market share was down by 7ppt to 20\% }
\end{figure}
\begin{figure}[htbp]
\centering
\includegraphics[width=0.88\linewidth]{figures/fig_117.jpg}
\caption{Figure 2 - Figure 2. A reconciliation between tepid EV sales vol and strong battery installation © 2026 Citi Inc. No redistribution without Citi's written permission. Figure 3. Who supplies to whom © 2026 Citi Inc. No redistribution without Citi's written permission.}
\end{figure}
\begin{figure}[htbp]
\centering
\includegraphics[width=0.88\linewidth]{figures/fig_155.jpg}
\caption{Exhibit 1 - Exhibit 3: We Estimate That Global Refinery Runs Declined 6.5mb/d Year-Over-Year at the End of July on Continuing Attacks on Refineries in the Middle East and Russia}
\end{figure}
\begin{figure}[htbp]
\centering
\includegraphics[width=0.88\linewidth]{figures/fig_158.jpg}
\caption{Exhibit 3 - Exhibit 5: Valuations of E\&P stocks reflect WTI oil prices slightly below the five-year strip Historical WTI price implied in covered E\&Ps and forward oil futures, \textbackslash{}\$/bbl}
\end{figure}
{\small\color{refgray}本节综合 6 篇研究；涉及 JPM、MS、Citi、GS。}
\newpage
\subsection{紫金黄金：终止收购后的战略重估与增长路径}
\textbf{紫金黄金终止收购Allied Gold并转为9.2\%少数股权认购，消除了压制估值的并购不确定性，市场焦点转向内生增长与潜在重启收购的期权价值。}
\subsubsection*{主流共识}
\begin{itemize}
  \item 终止收购Allied Gold消除了自5月以来压制紫金黄金股价的不确定性，市场情绪趋于缓和。
  \item 紫金黄金2026-2028年产量增长将主要依赖自有矿山爬坡，2030年100吨长期目标不变。
  \item 认购9.2\%股权保留了未来合作或重启收购的灵活性，体现了分步决策的审慎策略。
\end{itemize}
\subsubsection*{分歧与条件差异}
\begin{itemize}
  \item JPM认为终止收购是估值障碍的移除，而Citi更强调股权认购为未来重启收购保留了期权价值。
  \item JPM下调盈利预测12\%-32\%，但维持长期增长逻辑；Citi则关注短期业绩弹性（上半年利润同比增70\%+）和锂产能贡献。
  \item 对估值锚点的判断不同：JPM以FY27E P/E约13倍衡量，Citi采用DCF（WACC 8.2\%，终端增长率2.5\%）。
\end{itemize}
\subsubsection*{机构观点}
\paragraph{JPM} 终止收购Allied Gold消除了自5月以来压制股价的不确定性，紫金黄金相对同业落后约20\%的折价有望修复。剔除Allied并表贡献后，2026-2028年黄金产量预期调整为59/64/73吨，2028年70-75吨目标不变，Porgera等现有矿山爬坡是主要支撑。即便按当前金价重估，FY27E市盈率约13倍，估值仍具吸引力。盈利预测下调12\%-32\%，但2030年100吨长期产量目标未变；金价回调可能创造新的并购机会。
{\small\color{refgray}数据：紫金黄金相对黄金同业落后约20\%\par}
{\small\color{refgray}数据：FY27E市盈率约13倍\par}
{\small\color{refgray}数据：2026-2028年产量预期59/64/73吨\par}
{\small\color{refgray}数据：盈利预测下调12\%-32\%\par}
{\small\color{refgray}边际变化：从并购驱动的增长叙事转向内生增长，估值锚点从并表预期转为自有矿山爬坡。\par}
\paragraph{Citi} 紫金矿业终止对Allied Gold的全资收购（原方案每股44加元），改为以每股32.55加元认购1280万股，总对价约4.17亿加元，完成后持股9.2\%，折让约26\%。这一结构变化将一次性决策拆分为分步决策，保留了未来合作空间；若马里地缘政治不确定性下降，公司仍可能重启收购。预计上半年利润同比增长70\%以上，12万吨锂（LCE）产能计划将贡献下半年利润。股价自高点回调超30\%，估值吸引力显现，维持首选股。
{\small\color{refgray}数据：认购价32.55加元/股，较原方案折让26\%\par}
{\small\color{refgray}数据：总对价4.17亿加元（约2.95亿美元）\par}
{\small\color{refgray}数据：持股9.2\%\par}
{\small\color{refgray}数据：上半年利润同比增长70\%以上\par}
{\small\color{refgray}边际变化：从全资收购转向少数股权认购，强调分步决策和未来重启收购的期权价值。\par}
\subsubsection*{关键数据}
\begin{itemize}
  \item 紫金黄金相对黄金同业落后约20\%
  \item FY27E市盈率约13倍
  \item 2026-2028年产量预期59/64/73吨
  \item 盈利预测下调12\%-32\%
  \item 认购价32.55加元/股，较原方案折让26\%
  \item 总对价4.17亿加元（约2.95亿美元）
  \item 持股9.2\%
  \item 上半年利润同比增长70\%以上
  \item 12万吨锂（LCE）产能计划
  \item 股价自高点回调超30\%
\end{itemize}
\begin{figure}[htbp]
\centering
\includegraphics[width=0.88\linewidth]{figures/fig_080.jpg}
\caption{Figure 1 - Figure 1: Zijin Gold: Gold Production Targets Price Performance — 2259.HK Price (HK\textbackslash{}\$) — HSI (rebased)}
\end{figure}
\begin{figure}[htbp]
\centering
\includegraphics[width=0.88\linewidth]{figures/fig_081.jpg}
\caption{Figure 1 - Figure 1: Zijin Gold: Gold Production Targets Price Performance — 2259.HK Price (HK\textbackslash{}\$) — HSI (rebased) Company Data}
\end{figure}
{\small\color{refgray}本节综合 2 篇研究；涉及 JPM、Citi。}
\newpage
\subsection{价格战与产能再配置：中国增程SUV定价冲击与日本供应链转型}
\textbf{小米SkyNomad以低于同级竞品16-18万元的预售价切入增程SUV市场，可能重塑30万元级家庭SUV价格锚点；比亚迪7月基本面企稳但资金流仍在调整，8月业绩是关键验证点；松下将EV电池产线转向AI数据中心超级电容，显示日本供应链在EV需求放缓下的资产再配置。}
\subsubsection*{主流共识}
\begin{itemize}
  \item 小米SkyNomad N90/N70预售价（29.99万/25.99万元）显著低于理想L9和问界M9，形成明显价格优势，可能冲击中国高端家庭SUV市场格局。
  \item 比亚迪7月订单和批发数据企稳，库存环比下降6.0\%为年内最大降幅，渠道去库存加速，为下半年销售旺季腾出空间。
  \item 增程SUV市场是当前中国车市增长最快的细分领域之一，2025年PHEV\&EREV SUV销量达290万辆，小米切入该赛道具备市场基础。
  \item 松下将大阪EV电池产线转为超级电容生产，用于AI数据中心CBU，体现集团层面资产协同和业务重心向AI基础设施转移。
\end{itemize}
\subsubsection*{分歧与条件差异}
\begin{itemize}
  \item 小米SkyNomad的销量预期存在分歧：Bernstein预计N90月交付超1万台、N70约1.5万台，而GS基础预测2026年11万台、2027年24万台，乐观情形下2027年可达50万台，差异源于对燃油车置换需求释放速度的判断。
  \item 比亚迪基本面是否真正筑底存在分歧：Citi认为7月数据企稳是积极信号，但资金流仍在调整（AI和商用车板块本月继续跌26\%），8月业绩若不能证明国内盈利，情绪拐点可能延后。
  \item 小米增程系统的技术优势是否可持续：GS强调昆仑平台原生设计和馈电油耗优势（比竞品低0.6-1.6L/100km），但Bernstein对座舱灵活性功能（如座椅旋转）的实用场景持保留态度，认为竞争核心仍是价格。
  \item 松下产线改造的协同效应是否可复制：GS认为跨公司资产调配体现集团协同，但指出客户可能引入其他供应商、AI相关BBU业务竞争加剧等风险，改造的实际经济性需待量产验证。
\end{itemize}
\subsubsection*{机构观点}
\paragraph{Bernstein} 小米SkyNomad N90/N70预售价分别定为29.99万和25.99万元，显著低于理想L9（45.98万起）和问界M9 EREV（47.98万起），价差达16-18万元，直接改变30万元级家庭SUV购买决策的价格锚点。基于价格优势和同级竞品销量基准，预计N90月交付可超1万台，N70约1.5万台，但需注意N70与N90之间可能存在的内部竞争。正式上市价通常比预售价再低1万元以上，若兑现将进一步增强吸引力。
{\small\color{refgray}数据：N90 Max预售价29.99万元，N70 Max预售价25.99万元\par}
{\small\color{refgray}数据：N90纯电续航464km，综合续航1705km\par}
{\small\color{refgray}数据：N70纯电续航最高505km\par}
{\small\color{refgray}数据：理想L9起步价45.98万元，问界M9 EREV起步价47.98万元\par}
{\small\color{refgray}边际变化：小米以低于同级竞品三成以上的价格切入增程SUV市场，改变了此前市场对小米汽车定位的预期，价格锚点效应成为核心竞争变量。\par}
\paragraph{Citi} 比亚迪7月订单环比下降5\%，但剔除6月高基数后基本稳定；周度批发数据从7月中旬的-13\%收窄至月末的-1\%，7月国内新能源零售同比-6.8\%（较6月-9.9\%收窄）。加权基本面指数回落至0的稳态水平，库存环比下降6.0\%为年内最大降幅，显示渠道去库存加速。8月二季度业绩（预计净利润90-100亿元）是验证国内盈利能力和份额恢复的关键节点，若国内不亏，2027年出口270万辆对应50亿以上净利润。人形机器人业务首次纳入长期盈利框架，零部件与算法平台与汽车业务有60-80\%重合度。
{\small\color{refgray}数据：7月订单环比-5\%，批发周度数据从-13\%收窄至-1\%\par}
{\small\color{refgray}数据：7月库存环比-6.0\%，为年内最大降幅\par}
{\small\color{refgray}数据：二季度净利润预计90-100亿元\par}
{\small\color{refgray}数据：2027年出口目标270万辆，对应净利润超50亿元\par}
{\small\color{refgray}边际变化：比亚迪基本面从6月高基数回落中企稳，但资金流仍在调整（AI和商用车板块本月跌26\%），8月业绩成为情绪拐点的关键验证。\par}
\paragraph{GS} 松下计划将大阪EV电池产线改造为超级电容产线，目标2027年3月财年内量产，用于AI数据中心电容备份单元（CBU）。此举体现集团层面跨公司资产协同（松下能源与松下工业分属不同运营主体），电极材料和卷绕工艺的共通性降低改造技术门槛，盘活现有厂房设备避免新建资本开支。超级电容在响应速度和循环寿命上优于锂电池BBU，与AI数据中心供电需求匹配。参考报价5000日元，基于FY3/28E EV/EBITDA 9倍。
{\small\color{refgray}数据：产线改造目标2027年3月财年内量产\par}
{\small\color{refgray}数据：大阪两座EV电池工厂具备改造条件\par}
{\small\color{refgray}数据：超级电容CBU用于AI数据中心断电保护\par}
{\small\color{refgray}数据：参考报价5000日元，基于FY3/28E EV/EBITDA 9倍\par}
{\small\color{refgray}边际变化：松下从EV电池产线转向AI数据中心超级电容，标志其业务重心从汽车电动化向AI基础设施供电迁移，集团协同效应开始显现。\par}
\paragraph{GS} 小米SkyNomad系列以25.99万/29.99万元预售价切入增程SUV市场，昆仑制造平台原生设计带来平地板和长滑轨系统，减少约200毫米垂直侵入，提供4.4平方米规整空间，友商短期难以复制。N90 Max CLTC综合续航1705公里，纯电续航505公里，馈电油耗比竞品低0.6-1.6L/100km，支持92号汽油，降低燃油车用户置换门槛。预售价具有破坏性，可能扰动200k-400k价格带市场格局。基础预测2026年11万台、2027年24万台，乐观情形下2027年可达50万台。
{\small\color{refgray}数据：N70 Max预售价25.99万元，N90 Max预售价29.99万元\par}
{\small\color{refgray}数据：N90 Max综合续航1705km，纯电续航505km\par}
{\small\color{refgray}数据：馈电油耗比竞品低0.6-1.6L/100km\par}
{\small\color{refgray}数据：昆仑平台减少约200mm垂直侵入，提供4.4平方米空间\par}
{\small\color{refgray}边际变化：小米以原生增程平台和低于竞品三成以上的定价进入市场，GS强调其技术壁垒和价格破坏性，市场关注点转向9月正式上市前的配置披露。\par}
\subsubsection*{关键数据}
\begin{itemize}
  \item 小米N90 Max预售价29.99万元，较理想L9起步价低16万元，较问界M9 EREV低18万元
  \item N90纯电续航464km，综合续航1705km；N70纯电续航最高505km
  \item Bernstein预计N90月交付超1万台，N70约1.5万台
  \item 比亚迪7月库存环比-6.0\%，为年内最大降幅；批发周度数据从-13\%收窄至-1\%
  \item 比亚迪二季度净利润预计90-100亿元，2027年出口目标270万辆对应净利润超50亿元
  \item 松下将大阪EV电池产线转为超级电容产线，目标2027年3月财年内量产，用于AI数据中心CBU
  \item GS参考报价5000日元，基于FY3/28E EV/EBITDA 9倍
  \item 小米昆仑平台减少约200mm垂直侵入，馈电油耗比竞品低0.6-1.6L/100km
\end{itemize}
\begin{figure}[htbp]
\centering
\includegraphics[width=0.88\linewidth]{figures/fig_043.jpg}
\caption{EXHIBIT 1 - EXHIBIT 2: Xiaomi SkyNomad N90 Max is a full size 7-seater EREV SUV, with a highly competitive pre-sale price tag at RMB 299.9k; based on the sales volume and MSRP of its peer group, we estimate N90 Max can deliver over 10k units m}
\end{figure}
\begin{figure}[htbp]
\centering
\includegraphics[width=0.88\linewidth]{figures/fig_118.jpg}
\caption{Figure 1 - Figure 1. BYD weighted-fundamental index © 2026 Citi Inc. No redistribution without Citi's written permission. Figure 2. BYD fundamentals and its H share price movement BYD weighted-fundamental index (LHS) BYD-H share price movement MoM (RHS) © 2026 Citi Inc. }
\end{figure}
\begin{figure}[htbp]
\centering
\includegraphics[width=0.88\linewidth]{figures/fig_191.jpg}
\caption{Exhibit 5 - Exhibit 6: Xiaomi SU7/YU7 models accounted for 15\%/12\% of China's sedan/SUV markets with starting price at Rmb200-300k in Jan-May 2026 Xiaomi SU7/YU7 shipment and volume share \#\# Investment Thesis - Xiaomi Corp. We see Xiaomi, wh}
\end{figure}
\begin{figure}[htbp]
\centering
\includegraphics[width=0.88\linewidth]{figures/fig_044.jpg}
\caption{EXHIBIT 1 - EXHIBIT 3: Xiaomi SkyNomad N70 Max is a mid-large size 5-seater EREV SUV, with a pre-sale price set at RMB 259.9k; we estimate it can potentially sell 15k units monthly, while noting there could be cannibalization with N90, dependi}
\end{figure}
{\small\color{refgray}本节综合 4 篇研究；涉及 Bernstein、Citi、GS。}
\newpage
\subsection{区域市场与主题：政策转向、产业重构与结构性机会}
\textbf{全球市场正处于政策与产业双重转折点：中国政策重心从刺激转向落实，日本面临日元升值与AI驱动的产业重估，美国机器人限制重塑供应链，而光通信与地产数据则显示结构性改善。投资者需关注政策执行节奏、汇率波动与供应链本地化带来的差异化机会。}
\subsubsection*{主流共识}
\begin{itemize}
  \item 中国下半年政策以加快落实已规划措施为主，财政发力优先于货币宽松，政府债券发行提速是核心抓手。
  \item 日本光通信行业受AI数据中心需求拉动，出口量价齐升，日本厂商估值低于美国同行，存在修复空间。
  \item 美国对中国机器人限制对整机厂商冲击大于零部件厂商，已布局海外产能的零部件企业可能受益。
  \item 中国核心城市房地产成交量企稳，二手房价格止跌，8-10月销售有望在低基数下延续回升。
\end{itemize}
\subsubsection*{分歧与条件差异}
\begin{itemize}
  \item 中国政策宽松力度：GS认为存量债券额度（1.8万亿元）可支撑财政发力，但地方执行摩擦可能拖累落地；NOM预计下半年净融资6.8万亿元，但规模有限，且更依赖财政而非货币。
  \item 日元升值路径：JPM提示8月为关键窗口，极端情景下日元或走强15日元，但基准情形为有序升值；市场对升值速度与幅度存在分歧。
  \item 机器人限制影响：Citi认为中国零部件厂商因已海外布局而受益，但整机厂商面临需求萎缩；市场对供应链重构的受益范围与持续性看法不一。
  \item 地产复苏持续性：Citi对8-10月销售回升持积极态度，但土地购置大幅下滑（-22\%）可能制约中期供给，复苏斜率存在不确定性。
\end{itemize}
\subsubsection*{机构观点}
\paragraph{JPM} 8月是日元阶段性走强的关键窗口，美元兑日元接近4Q26目标164，日本利率市场变化（日债收益率上升）使政策焦点从汇率稳定转向抑制长端利率。GPIF 8月7日持仓数据若显示国内债券比例上升，将强化资金回流预期。极端情景下日元或走强15日元，但基准为有序升值。在此环境下，强势日元篮子（94\%国内收入）和进口商篮子（78\%国内收入）相对稳健，而出口商和AI相关板块仓位拥挤，轮动风险大于指数涨跌。韩国回购篮子作为防御性分散配置，在波动市况中具有韧性。
{\small\color{refgray}数据：美元兑日元4Q26目标164\par}
{\small\color{refgray}数据：日元极端走强幅度15日元\par}
{\small\color{refgray}数据：强势日元篮子94\%收入来自国内\par}
{\small\color{refgray}数据：进口商篮子78\%收入来自国内\par}
{\small\color{refgray}边际变化：从单纯看跌日元转向关注日元升值风险，强调结构性轮动而非指数方向。\par}
\paragraph{NOM} 7月政治局会议重新强调“逆周期调节”，确认政策立场转向宽松，但发力规模有限，路径上更依赖财政而非货币。二季度GDP为疫情后最低，固投和零售销售同比分别降至-9.7\%和0.2\%，印证“两个K型”分化。下半年政府债券净融资预计6.8万亿元，占全年额度57\%，较去年同期5.75万亿元提速。会议首提“贸易再平衡”，取消部分出口退税、增加美农产品采购等已在推进，但政策支持力度受制于货币工具空间与效果。
{\small\color{refgray}数据：二季度固投同比-9.7\%\par}
{\small\color{refgray}数据：零售销售同比0.2\%\par}
{\small\color{refgray}数据：下半年政府债券净融资6.8万亿元\par}
{\small\color{refgray}数据：占全年额度57\%\par}
{\small\color{refgray}边际变化：政策措辞从“跨周期”转向“逆周期”，明确宽松信号，但强调财政为主、货币为辅。\par}
\paragraph{Citi} 美国FCC将人形和四足机器人列入Covered List，对中国整机厂商冲击最大，因其配备传感器收集敏感数据，监管审查更严。但已嵌入美国供应链的零部件公司反而可能获得更大美国市场份额，因美国整机厂需依赖中国成熟零部件体系。恒立液压将行星滚柱丝杠产线迁往墨西哥，绿的谐波与敏实在美成立执行器合资公司，荣泰和伟创在泰国设厂，这些提前布局使零部件商规避直接出口限制。
{\small\color{refgray}数据：恒立液压产线迁往墨西哥\par}
{\small\color{refgray}数据：绿的谐波在美合资公司\par}
{\small\color{refgray}数据：荣泰、伟创泰国工厂\par}
{\small\color{refgray}数据：禁令生效日期2026年7月28日\par}
{\small\color{refgray}边际变化：从关注整机受限转向挖掘零部件本地化受益标的，供应链地图成为关键。\par}
\paragraph{Citi} 7月34城新房成交同比转正至11\%，18城二手房同比+6.8\%，一线城市二手房同比+14.1\%，佛山、宁波等二线增速超20\%。6个核心城市二手房价格环比企稳，为新房定价提供参照。上半年5家房企销售正增长（中海+12\%、金茂+8\%等），预计8-9月更多房企转正。土地市场上半年300城成交建面同比-22\%，但优质地块溢价率升至13.7\%。政策端中央预算内投资970亿用于城市更新（+21\%），1600亿超长期特别国债支持。8-10月行业具备阶段性表现条件。
{\small\color{refgray}数据：34城新房成交同比+11\%\par}
{\small\color{refgray}数据：一线城市二手房同比+14.1\%\par}
{\small\color{refgray}数据：佛山、宁波二手房增速超20\%\par}
{\small\color{refgray}数据：上半年300城土地成交建面同比-22\%\par}
{\small\color{refgray}边际变化：从关注成交下滑转向确认企稳信号，强调低基数与推盘放量带来的阶段性机会。\par}
\paragraph{GS} 7月政治局会议确认政策重心为加快落实已规划措施，而非新刺激。财政方面，存量未使用政府债券额度约1.8万亿元（中央0.6万亿、地方1.2万亿）可较快调用，新增8000亿元政策性金融工具（高于去年5000亿）。货币方面，重申适度宽松，但删去“保持流动性合理充裕”，工具运用将更灵活。地方换届、反腐等不确定性可能影响资金落地速度。
{\small\color{refgray}数据：存量未使用政府债券额度1.8万亿元\par}
{\small\color{refgray}数据：中央0.6万亿元\par}
{\small\color{refgray}数据：地方1.2万亿元\par}
{\small\color{refgray}数据：新增政策性金融工具8000亿元\par}
{\small\color{refgray}边际变化：从预期新刺激转向关注存量工具加速落地，强调执行摩擦风险。\par}
\paragraph{GS} 日本6月光缆出口额同比+129\%，光纤+56\%，均价+55\%，对美出口额同比+50\%（环比+95\%），对美光纤出口额+109\%（环比+138\%）。对美光缆出口量-16\%但均价+78\%，显示高附加值产品占比提升。横滨辖区出口额+226\%（环比+419\%），对应住友电气出货加速。估值方面，康宁32倍PE，藤仓24倍、住友18倍、古河20倍，差距可能收窄。住友电气信息通信业务存在上调指引空间，藤仓下半年指引保守，古河液冷板下半年贡献业绩。
{\small\color{refgray}数据：光缆出口额同比+129\%\par}
{\small\color{refgray}数据：光纤出口额+56\%\par}
{\small\color{refgray}数据：光纤均价+55\%\par}
{\small\color{refgray}数据：对美光缆出口额+50\%\par}
{\small\color{refgray}边际变化：从关注AI需求转向验证日本厂商量价齐升与估值修复逻辑，海关数据提供微观证据。\par}
\subsubsection*{关键数据}
\begin{itemize}
  \item 中国二季度固投同比-9.7\%，零售销售0.2\%
  \item 下半年政府债券净融资6.8万亿元，占全年57\%
  \item 存量未使用政府债券额度1.8万亿元
  \item 日本6月光缆出口额同比+129\%，对美+50\%
  \item 康宁PE 32倍 vs 藤仓24倍、住友18倍
  \item 34城新房成交同比+11\%，一线二手房+14.1\%
  \item 上半年300城土地成交建面同比-22\%
  \item 美元兑日元4Q26目标164，极端升值15日元
\end{itemize}
\begin{figure}[htbp]
\centering
\includegraphics[width=0.88\linewidth]{figures/fig_061.jpg}
\caption{Figure 1 - Figure 1: Strong Yen and Importers have worked best in orderly USDJPY declines, while in fast de-risking episodes their value has been mainly relative Figure 2: The Korea buyback basket has shown resilience during the recent mom}
\end{figure}
\begin{figure}[htbp]
\centering
\includegraphics[width=0.88\linewidth]{figures/fig_063.jpg}
\caption{Figure 3 - Figure 3: With USDJPY approaching our house 4Q26 target of 164, August looks like a key window for episodic yen-strength risk Figure 4: With JGB yields rising, market focus has shifted from currency stabilization to policies aim}
\end{figure}
\begin{figure}[htbp]
\centering
\includegraphics[width=0.88\linewidth]{figures/fig_172.jpg}
\caption{Exhibit 7 - Exhibit 1: Export value for optical fiber cable (¥ mn) Exhibit 2: Export value for optical fiber (¥ mn) Includes some other optical fiber cables. Exhibit 3: Export volume for optical fiber cable (kg)}
\end{figure}
\begin{figure}[htbp]
\centering
\includegraphics[width=0.88\linewidth]{figures/fig_174.jpg}
\caption{Exhibit 2 - Exhibit 5: Average export price for optical fiber cable (¥1,000/kg)}
\end{figure}
{\small\color{refgray}本节综合 6 篇研究；涉及 JPM、NOM、Citi、GS。}
\newpage
\subsection{AI算力扩张重塑半导体与硬件价值链}
\textbf{AI算力需求正从云端向存储、光通信、电源管理及终端设备全链条传导，驱动定价权转移与产品结构升级。光刻强度提升、HDD定价加速、AI电源需求爆发是三大核心主线，但各环节盈利兑现节奏与可持续性存在显著分化。}
\subsubsection*{主流共识}
\begin{itemize}
  \item AI数据中心资本开支扩张是半导体设备、存储、光通信等环节需求增长的核心驱动力，订单能见度普遍延伸至2027年以后
  \item 存储行业供需缺口持续扩大，HDD与NAND定价权显著增强，供应商连续多个季度实现超预期业绩并上调指引
  \item 光刻强度（litho intensity）在DRAM和逻辑芯片中持续提升，ASML等设备商受益于EUV/DUV双重需求增长
  \item AI电源管理（含UPS/BBU电池、功率半导体）成为新的高增长细分市场，相关企业收入指引普遍超过70\%同比增长
\end{itemize}
\subsubsection*{分歧与条件差异}
\begin{itemize}
  \item AI对利润率的影响方向存在分歧：UBS认为AWS利润率可维持38\%水平，而MS指出高通AI业务初期毛利率低于平均水平，显示AI收入质量因环节而异
  \item 存储涨价持续性判断不一：JPM认为HDD定价增速将从11\%跳升至20\%以上，而GS认为存储成本上升将压制中国智能手机终端需求，全年出货量预计下降10\%
  \item 苹果供应链调整节奏存在预期差：MS认为高通在iPhone 18份额低于20\%假设，而JPM观察到Skyworks苹果业务内容下滑幅度收窄至低两位数，显示不同供应商受影响程度不同
  \item 对ASML增长驱动力的判断分化：Bernstein强调DRAM取代逻辑成为主要引擎，而部分市场观点仍聚焦于逻辑先进制程扩产
\end{itemize}
\subsubsection*{机构观点}
\paragraph{BARC} 开云集团信用状况显著改善，净债务从80亿欧元降至33亿欧元，杠杆倍数从3.4倍降至1.4倍。Gucci二季度有机收入下滑收窄至-2\%，好于市场预期的-4.7\%，皮具品类恢复增长，营业利润率同比提升100个基点至17\%。出售Kering Beauté回笼40亿欧元现金，叠加Valentino期权延后至2028-29年，为管理层聚焦品牌修复争取时间。预计EBIT利润率到2029年将达22.2\%，较资本市场日目标提前一年。
{\small\color{refgray}数据：净债务80亿→33亿欧元\par}
{\small\color{refgray}数据：Gucci有机收入-2\% vs 预期-4.7\%\par}
{\small\color{refgray}数据：Gucci营业利润率17\%（+100bps）\par}
{\small\color{refgray}数据：EBIT利润率2029年目标22.2\%\par}
{\small\color{refgray}边际变化：新CEO上任后快速去杠杆，出售非核心资产回笼现金，评级展望由负面调回稳定\par}
\paragraph{Bernstein} ASML光刻强度正从24\%向2028年28\%持续抬升，定价能力尚未被市场充分计入。DRAM取代逻辑成为光刻需求主要引擎，存储系统销售占比从30\%升至50\%，2026年存储销售预计增长超75\%。DRAM EUV曝光量未来五年将增长4.5倍，从340万片/月升至2630万片/月。EUV从3800E向3800F迁移带来约20\%价格上浮，DUV在1:1配套比例下构成结构性增长来源。
{\small\color{refgray}数据：光刻强度24\%→28\%（2028年）\par}
{\small\color{refgray}数据：存储销售占比30\%→50\%\par}
{\small\color{refgray}数据：DRAM EUV曝光量5年增长4.5倍\par}
{\small\color{refgray}数据：EUV 3800F价格上浮约20\%\par}
{\small\color{refgray}边际变化：增长逻辑从逻辑芯片单一驱动转向DRAM+逻辑双引擎，HNA和双重曝光EUV将DRAM光刻强度推至30\%\par}
\paragraph{JPM} 英飞凌3Q26业绩大概率符合预期，核心看点在AI电源需求与连续提价带来的指引上修。汽车半导体回暖获同行验证，TI汽车业务同比增长中双位数，NXP增长17\%且为内容驱动。4月和7月连续两次提价，叠加数据中心CPU功率半导体需求强劲，公司具备上调收入和利润率指引条件。AI电源市场FY28规模可达140亿欧元，英飞凌按28.5\%份额对应约40亿欧元收入。
{\small\color{refgray}数据：AI电源市场FY28达140亿欧元\par}
{\small\color{refgray}数据：英飞凌份额28.5\%对应40亿欧元\par}
{\small\color{refgray}数据：NXP汽车业务+17\%\par}
{\small\color{refgray}数据：4月和7月连续两次提价\par}
{\small\color{refgray}边际变化：汽车半导体需求定性为'周期开始'而非补库存，AI电源成为长期增长主逻辑\par}
\paragraph{JPM} 希捷F4Q26数据中心收入29.33亿美元，同比增长57\%，每EB定价同比增长11\%，下季度指引至少20\%。供需缺口持续扩大，近线产能已锁定至2028年，部分客户规划2029年需求。增量毛利率83\%以上，下季度指引88\%以上。NAND大幅涨价带来替代效应，Edge IoT获得机遇性定价。FY27和FY28每股盈利预测分别上调30.7\%和32.1\%。
{\small\color{refgray}数据：数据中心收入+57\%至29.33亿美元\par}
{\small\color{refgray}数据：每EB定价+11\%，下季指引≥20\%\par}
{\small\color{refgray}数据：增量毛利率83\%→88\%\par}
{\small\color{refgray}数据：FY27/FY28 EPS上调30.7\%/32.1\%\par}
{\small\color{refgray}边际变化：定价增速从4\%加速至11\%再跳升至20\%指引，连续第二个季度beat-and-raise\par}
\paragraph{JPM} Skyworks六月季度收入9.35亿美元超预期，EPS 1.08美元超共识。公司取消4.3亿美元季度股息，推出20亿美元回购，资金使用顺序为回购、去杠杆、并购。苹果业务内容下滑从20-25\%收窄至低两位数，驱动因素为客户出货量增长20-22\%及机型组合改善。射频复杂度多年来首次上升，但需下一次选型周期验证。光通信增速从30\%降至15\%，属基数效应。
{\small\color{refgray}数据：收入9.35亿美元 vs 预期9.25亿\par}
{\small\color{refgray}数据：EPS 1.08美元 vs 共识1.03\par}
{\small\color{refgray}数据：苹果内容下滑收窄至低两位数\par}
{\small\color{refgray}数据：20亿美元回购替代4.3亿股息\par}
{\small\color{refgray}边际变化：资本回报框架从股息转向回购，苹果内容下滑斜率明显放缓\par}
\paragraph{MS} 高通在iPhone 18基带份额低于20\%假设，苹果不再采用毫米波版本导致订单减少。但QTL专利授权收入保持完整，冲击集中在芯片销售而非专利费。AI业务12月季度爬坡至约6亿美元，但初期毛利率低于平均水平，本季度非GAAP毛利率54.0\%低于市场预期55.6\%。汽车业务年化收入从60亿上调至70亿美元，成为稳健增长引擎。FY27的50亿美元AI目标将改变叙事。
{\small\color{refgray}数据：iPhone 18份额<20\%\par}
{\small\color{refgray}数据：AI业务12月季度约6亿美元\par}
{\small\color{refgray}数据：毛利率54.0\% vs 预期55.6\%\par}
{\small\color{refgray}数据：汽车业务年化收入60→70亿美元\par}
{\small\color{refgray}边际变化：苹果退出节奏加快但QTL稳定，AI收入起量快于预期但毛利率稀释明显\par}
\paragraph{UBS} AWS 2Q26同比增长37\%，高于预期32\%。AI年化收入从100亿增至250亿美元，总订单积压4960亿美元，环比增加510亿。剔除6亿美元能源成本收益后，AWS运营利润率环比持平在38\%，AI计算负载尚未拉低利润率。预计2026下半年增速达39\%，明年可能突破40\%。资本开支上调至220亿美元，反映AI基础设施投入持续加码。
{\small\color{refgray}数据：AWS收入+37\% vs 预期32\%\par}
{\small\color{refgray}数据：AI年化收入100→250亿美元\par}
{\small\color{refgray}数据：订单积压4960亿美元（+510亿）\par}
{\small\color{refgray}数据：运营利润率持平38\%\par}
{\small\color{refgray}边际变化：AI订单积压规模超预期，市场对AI负载压缩利润率的担忧暂未验证\par}
\paragraph{GS} 中国智能手机6月出货1700万部，同比下降17\%，环比下降36\%，存储成本上升压制终端需求，全年出货量预期调整为同比下降10\%。二季度出货6900万部高于预期5900万部。单机摄像头数量从2022年峰值3.8颗降至2.9颗，但20MPx以上渗透率从2024年52\%升至2026年66\%，配置从数量堆叠转向像素升级。
{\small\color{refgray}数据：6月出货1700万部（同比-17\%）\par}
{\small\color{refgray}数据：全年出货预期同比-10\%\par}
{\small\color{refgray}数据：单机摄像头3.8→2.9颗\par}
{\small\color{refgray}数据：20MPx+渗透率52\%→66\%\par}
{\small\color{refgray}边际变化：存储成本上升成为终端需求核心变量，摄像头配置从数量转向像素升级\par}
\paragraph{GS} 三星SDI 2Q26恢复盈利，营业利润2040亿韩元，但剔除关税退还后核心亏损约130亿韩元。UPS和BBU电池2026年收入指引均同比增长超70\%，市场份额40-50\%，高功率圆柱产线满负荷运转。美国ESS订单已覆盖至2029年，2028年起需求将超现有产能。欧洲电动车需求现恢复迹象，xEV亏损预计收窄。
{\small\color{refgray}数据：营业利润2040亿韩元（含关税退还）\par}
{\small\color{refgray}数据：UPS/BBU收入指引+70\%\par}
{\small\color{refgray}数据：市场份额40-50\%\par}
{\small\color{refgray}数据：美国ESS订单覆盖至2029年\par}
{\small\color{refgray}边际变化：AI数据中心备用电源成为核心增长引擎，订单能见度从2-3年延伸至2029年\par}
\paragraph{GS} 索尼拟以每股1300日元收购腾龙剩余股份，溢价约15\%，总对价约1880亿日元。索尼已持有15.35\%股份，交易定位为'创作技术'领域战略投资，不改变各业务优先级。影像业务是ET\&S板块盈利核心，镜头是价值链关键环节。腾龙在普通用户市场产品矩阵完整，并向车载和医疗镜头拓展，与索尼车载图像传感器业务形成协同。
{\small\color{refgray}数据：报价1300日元/股，溢价15\%\par}
{\small\color{refgray}数据：总对价约1880亿日元\par}
{\small\color{refgray}数据：索尼持股15.35\%\par}
{\small\color{refgray}数据：腾龙市值约1830亿日元\par}
{\small\color{refgray}边际变化：索尼战略投资从三大娱乐业务向影像硬件延伸，补强供应链关键环节\par}
\paragraph{GS} VPEC向Aixtron采购6.53亿新台币MOCVD设备，交货周期12个月，订单能见度延伸至2027年后。3Q26营收预计环比增长23\%至11.86亿新台币，光电子业务占比从27\%跳升至41\%。客户已在InP磊晶片针对300mW CW激光器下达研发订单，CPO方向机会开始转化。上调2027-28年净利润预测2-3\%，毛利率预计从38.6\%升至3Q26的41\%和4Q26的45.7\%。
{\small\color{refgray}数据：MOCVD设备采购6.53亿新台币\par}
{\small\color{refgray}数据：3Q26营收环比+23\%\par}
{\small\color{refgray}数据：光电子占比27\%→41\%\par}
{\small\color{refgray}数据：毛利率38.6\%→41\%→45.7\%\par}
{\small\color{refgray}边际变化：AI光模块从800G向1.6T升级，InP衬底供应改善推动订单转化为出货\par}
\subsubsection*{关键数据}
\begin{itemize}
  \item ASML光刻强度从24\%升至2028年28\%，DRAM EUV曝光量5年增长4.5倍
  \item 希捷每EB定价增速从11\%跳升至20\%指引，增量毛利率88\%以上
  \item AWS AI年化收入从100亿增至250亿美元，订单积压4960亿美元
  \item 英飞凌AI电源市场FY28达140亿欧元，4月和7月连续提价
  \item 三星SDI UPS/BBU电池收入指引同比增长超70\%，市场份额40-50\%
  \item VPEC 3Q26光电子收入占比从27\%跳升至41\%，毛利率预计升至45.7\%
  \item 中国智能手机6月出货同比下降17\%，全年预期下调至-10\%
  \item 高通AI业务12月季度爬坡至6亿美元，但毛利率54.0\%低于预期55.6\%
\end{itemize}
\begin{figure}[htbp]
\centering
\includegraphics[width=0.88\linewidth]{figures/fig_026.jpg}
\caption{EXHIBIT 1 - EXHIBIT 2: We anticipate an increase in lithography intensity in DRAM. The introduction of double-patterning EUV (or HNA) at the 1d node is expected to drive lithography intensity up to \textbackslash{}\textasciitilde{}30\% EXHIBIT 3: Logic litho intensity decl}
\end{figure}
\begin{figure}[htbp]
\centering
\includegraphics[width=0.88\linewidth]{figures/fig_029.jpg}
\caption{Exhibit 5 - EXHIBIT 5: We estimate DRAM total EUV exposures to increase 4.5x over the next 5 years, increasing from 26\% to 44\% of total exposures. While we expect logic EUV exposure to grow at a slower pace than DRAM, we still expect it to mor}
\end{figure}
\begin{figure}[htbp]
\centering
\includegraphics[width=0.88\linewidth]{figures/fig_098.jpg}
\caption{Figure 1 - Figure 1: Amazon.com, Inc. - Unit Growth vs Shipping Cost Growth (1Q18-2Q26) Reported net sales of \textbackslash{}\$200.6 billion vs guidance range of \textbackslash{}\$194 billion to \textbackslash{}\$199 billion (+16\%-19\% YOY) came in above the Street \textbackslash{}\$196.8 billion and U}
\end{figure}
\begin{figure}[htbp]
\centering
\includegraphics[width=0.88\linewidth]{figures/fig_188.jpg}
\caption{Exhibit 1 - Exhibit 1: VPEC revenues by products and GM trend 3-month revenue preview: VPEC delivered flat revenues QoQ in 1Q26 and 2Q26; however, we expect +23\% QoQ in 3Q26 thanks to capacity expansion and better InP substrate supply, on th}
\end{figure}
{\small\color{refgray}本节综合 11 篇研究；涉及 BARC、Bernstein、JPM、MS、UBS、GS。}
\newpage
\subsection{股票策略、估值与资金流：业绩可靠性重估与长期指引上修}
\textbf{本期机构观点聚焦于个股层面盈利质量与长期增长可见度的重新定价：Bernstein认为Air Products已通过'鸭子测试'，应回归工业气体估值框架；GS上调Equinix长期增长指引并维持估值倍数，同时上调赛诺菲短期盈利但下调长期估值，反映市场对管线接续能力的担忧。整体看，市场正从短期业绩超预期转向对可持续增长路径的审视。}
\subsubsection*{主流共识}
\begin{itemize}
  \item 连续超预期并上调指引（beat-and-raise）是重建业绩可靠性的关键，也是估值溢价的基础。
  \item 长期增长可见度（如多年指引、项目时间表）对估值的影响大于单季业绩。
  \item 资本开支扩张与回报率假设的稳定性是支撑长期估值的重要条件。
  \item 研发管线或项目融资的不确定性仍是估值折价的核心来源。
\end{itemize}
\subsubsection*{分歧与条件差异}
\begin{itemize}
  \item Air Products的NEOM项目是估值拖累还是已被过度惩罚？Bernstein认为影响可控，市场可能误读。
  \item Equinix大幅上调资本开支但维持回报率假设，市场对其执行风险存在分歧。
  \item 赛诺菲短期盈利上修与长期估值下修并存，反映市场对管线接续能力的不同预期。
  \item 对Air Products FY27增速放缓的解读：是短期波动还是长期趋势？
\end{itemize}
\subsubsection*{机构观点}
\paragraph{Bernstein} Bernstein认为Air Products已通过'鸭子测试'，其收入与经营表现符合工业气体公司特征，应据此定价。连续超预期并上调指引重建了业绩可靠性，NEOM项目时间表调整后对FY27无重大财务影响，拆除了估值折价的核心论据。FY26调整后EPS增速约11-12\%，FY27指引6-9\%属保守，长期复利逻辑不变。上调FY27 EPS至14.67美元，维持26倍NTM市盈率，对应约27\%上行空间。
{\small\color{refgray}数据：FY26调整后EPS增速11-12\%\par}
{\small\color{refgray}数据：FY27 EPS预测上调至14.67美元\par}
{\small\color{refgray}数据：26倍NTM市盈率\par}
{\small\color{refgray}数据：目标价373美元，较发布日收盘价约27\%空间\par}
{\small\color{refgray}边际变化：从'NEOM拖累估值'转向'NEOM影响可控'，并首次明确按工业气体框架定价。\par}
\paragraph{GS} GS认为Equinix二季度超预期且长期指引大幅上修，核心是AI推理负载早期增长与零售托管需求强劲。资本开支区间从30-40亿上调至50-70亿美元（85\%用于扩张），同时维持25\%现金回报率假设，显示管理层对需求可见度的信心。2027-29年收入增速指引从7-10\%上调至10-13\%，AFFO每股增速从5-9\%上调至9-12\%，2029年AFFO每股约58美元。估值倍数维持21倍，变化完全来自盈利上修。
{\small\color{refgray}数据：2Q26调整后EBITDA 13.96亿美元，高于预期13.65亿\par}
{\small\color{refgray}数据：AFFO每股11.78美元，超预期11.46美元\par}
{\small\color{refgray}数据：资本开支上调至50-70亿美元，85\%用于扩张\par}
{\small\color{refgray}数据：2029年AFFO每股58美元，共识56.47美元\par}
{\small\color{refgray}边际变化：从'短期超预期'转向'长期增长可见度提升'，资本开支大幅上调但回报率假设不变。\par}
\paragraph{GS} GS对赛诺菲采取短期盈利上修与长期估值下修的组合：2026-27年EPS上调2-3\%至8.56和9.72欧元，反映汇率、Dupixent强劲增长及疫苗节奏；但2028年后因研发投入增加和终止资产不再贡献收入，EPS下调约1\%。目标市盈率从8.5倍降至8倍，DCF估值从83欧元降至75欧元。财报后股价跌9\%，GS认为反应过度，但管线接续能力仍是估值折价根源。
{\small\color{refgray}数据：2026年EPS上调至8.56欧元，2027年9.72欧元\par}
{\small\color{refgray}数据：目标市盈率从8.5倍下调至8倍\par}
{\small\color{refgray}数据：DCF估值从83欧元降至75欧元\par}
{\small\color{refgray}数据：财报后股价下跌约9\%\par}
{\small\color{refgray}边际变化：从'商业执行强劲'转向'短期盈利上修与长期估值下修并存'，强调研发管线不确定性。\par}
\subsubsection*{关键数据}
\begin{itemize}
  \item Air Products FY27 EPS预测14.67美元，26倍NTM市盈率，目标价373美元
  \item Equinix 2Q26 AFFO每股11.78美元，超预期2.8\%
  \item Equinix 2027-29年收入增速指引上调至10-13\%，资本开支50-70亿美元
  \item 赛诺菲2026年EPS上调至8.56欧元，目标市盈率下调至8倍
  \item 赛诺菲财报后股价跌9\%，DCF估值降至75欧元
\end{itemize}
\begin{figure}[htbp]
\centering
\includegraphics[width=0.88\linewidth]{figures/fig_021.jpg}
\caption{EXHIBIT 1 - EXHIBIT 1: Air Products are starting to build the results reliability investors prize from the industrial gases APD Adj. EPS beat / miss history under Eduardo Menezes}
\end{figure}
\begin{figure}[htbp]
\centering
\includegraphics[width=0.88\linewidth]{figures/fig_187.jpg}
\caption{Exhibit 3 - Exhibit 3: GSe vs consensus \#\# Valuation and risks We are Neutral rated. We derive our EU large cap pharma price targets using a 50:50 blend of a DCF-based valuation and a P/E multiple-based approach. Our bottom-up DCF methodolog}
\end{figure}
\begin{figure}[htbp]
\centering
\includegraphics[width=0.88\linewidth]{figures/fig_022.jpg}
\caption{EXHIBIT 1 - EXHIBIT 3: Air Products' merchant pricing continues to keep pace with peers despite helium headwinds}
\end{figure}
\begin{figure}[htbp]
\centering
\includegraphics[width=0.88\linewidth]{figures/fig_023.jpg}
\caption{EXHIBIT 1 - EXHIBIT 3: Air Products' merchant pricing continues to keep pace with peers despite helium headwinds EXHIBIT 4: Recent margin progress is visible across all regions}
\end{figure}
{\small\color{refgray}本节综合 3 篇研究；涉及 Bernstein、GS。}
\newpage
\subsection{外汇与跨境资金：AI光通信分化与药企管线重估下的跨境资金流向}
\textbf{本期外汇与跨境资金栏目聚焦AI基础设施与医药研发两大主题的跨境资金动向：JPM认为AI光通信需求弹性高但运营商业务拖累整体，估值倍数下修；GS指出赛诺菲短期盈利扎实但管线缺口引发市场重定价，目标价下调。两者均显示跨境资金正从传统电信和成熟药企流向AI相关光通信等高弹性领域。}
\subsubsection*{主流共识}
\begin{itemize}
  \item AI相关光通信需求保持高弹性，订单增速加快，是跨境资金追逐的核心方向。
  \item 运营商级光通信业务因客户下单节奏波动，短期表现低于预期，但长期逻辑未变。
  \item 药企研发转型周期长，市场对管线接续核心产品的能力持谨慎态度。
  \item 估值倍数向行业同业靠拢成为趋势，JPM下调康宁至30倍，GS下调赛诺菲目标价。
\end{itemize}
\subsubsection*{分歧与条件差异}
\begin{itemize}
  \item JPM认为康宁业绩超卖方预期但未达买方标准，而GS认为赛诺菲股价反应过度，市场预期分歧明显。
  \item JPM上调康宁盈利预测，GS上调赛诺菲短期EPS但下调长期EPS，反映对增长持续性的不同判断。
  \item JPM强调光通信长期逻辑未变，GS则强调研发转型需七年以上，时间成本差异显著。
  \item 日元假设调整（JPM从120至150）与汇率因素（GS受益）对盈利影响方向相反。
\end{itemize}
\subsubsection*{机构观点}
\paragraph{JPM} 康宁2Q26业绩超卖方预期但未达买方标准，光通信内部分化：企业级AI相关销售弹性高、订单增速加快，运营商级业务因客户下单节奏波动低于预期。短期噪音不改长期AI光通信叙事，但估值倍数从35倍降至30倍，更贴近光通信同业水平，同时反映其美国本土垂直整合供应商的独特地位。目标价从200美元下调至170美元，基于过去12个月股价上涨127.4\%后的估值修正。
{\small\color{refgray}数据：2Q26收入47.38亿美元，高于JPM预期46.41亿和公司指引46亿\par}
{\small\color{refgray}数据：毛利率39.6\%，超出预期38.6\%\par}
{\small\color{refgray}数据：EPS 0.78美元，高于预期0.77美元\par}
{\small\color{refgray}数据：估值倍数从35倍降至30倍\par}
{\small\color{refgray}边际变化：估值倍数下调反映市场从AI光通信高增长预期转向更务实的同业比较，同时日元假设从120调整至150，对净利润产生负面影响但公司计划通过套保、定价和降本对冲。\par}
\paragraph{GS} 赛诺菲2Q26业绩后股价单日回落约9\%，短期盈利动能扎实，Dupixent商业执行超预期，GS上调2026-2027年EPS预测2-3\%；但amlitelimab、itepekimab等多项资产终止带来管线缺口，长期EPS下调约1\%，目标价从84欧元调整至79欧元。GS认为市场反应可能过度，但管理层需提供结构化方案重建信任，研发转型周期可能超过七年。
{\small\color{refgray}数据：股价单日回落约9\%\par}
{\small\color{refgray}数据：FY26/FY27 EPS预测上调2-3\%\par}
{\small\color{refgray}数据：长期EPS预测下调约1\%\par}
{\small\color{refgray}数据：目标价从84欧元调整至79欧元，隐含约9\%上行空间\par}
{\small\color{refgray}边际变化：市场关注点从短期业绩转向长期管线接续能力，研发失败与资产终止成为重新定价的核心变量，管理层可能于2026年下半年举办研发更新活动。\par}
\subsubsection*{关键数据}
\begin{itemize}
  \item 康宁2Q26收入47.38亿美元，超预期；毛利率39.6\%，EPS 0.78美元
  \item 康宁估值倍数从35倍降至30倍，目标价从200降至170美元
  \item 康宁过去12个月股价上涨127.4\%
  \item 康宁2026年资本开支约20亿美元，2027年约25亿美元
  \item 赛诺菲股价单日回落约9\%，目标价从84欧元调整至79欧元
  \item 赛诺菲FY26/FY27 EPS预测上调2-3\%，长期EPS下调约1\%
  \item 赛诺菲研发转型周期可能超过七年
  \item 日元假设从120调整至150，对康宁净利润产生负面影响
\end{itemize}
\begin{figure}[htbp]
\centering
\includegraphics[width=0.88\linewidth]{figures/fig_185.jpg}
\caption{Exhibit 3 - Exhibit 3: GSe vs consensus \#\# Valuation and risks We are Neutral rated. We derive our EU large cap pharma price targets using a 50:50 blend of a DCF-based valuation and a P/E multiple-based approach. Our bottom-up DCF methodolog}
\end{figure}
{\small\color{refgray}本节综合 2 篇研究；涉及 JPM、GS。}
\newpage
\subsection{能源、大宗与公用事业：LNG资产变现窗口与资本开支再平衡}
\textbf{在LNG价格高位、供需偏紧的背景下，中国石油考虑出售LNG Canada部分股权以筹措二期扩建资金，反映国际油气巨头正利用景气窗口优化资产组合、分担资本开支压力。市场对潜在减值影响存在分歧，但基本面强劲（二季度净利预计同比+56\%）提供支撑。}
\subsubsection*{主流共识}
\begin{itemize}
  \item LNG市场供需偏紧，价格处于高位，为资产变现提供有利估值窗口。
  \item LNG Canada二期项目资本开支规模大（约136.5亿美元），需引入合作方分担资金压力。
  \item 中国石油出售部分股权主要逻辑为回笼资金、降低二期资本开支负担，符合商业理性。
  \item 参考Petronas近期交易（30亿美元出售5\%股权及上游资产），当前时点出售可获得合理估值。
\end{itemize}
\subsubsection*{分歧与条件差异}
\begin{itemize}
  \item 市场担忧资产减值风险，但JPM认为是否计提取决于出售估值与账面价值对比，目前中性偏正面。
  \item 二期项目FID时间表存在不确定性（2026年底前），投产可能推迟至2030年代初，影响长期供应预期。
  \item 中国石油最新预测已保守计入150亿人民币减值，但实际减值规模可能因最终交易价格而偏离。
\end{itemize}
\subsubsection*{机构观点}
\paragraph{JPM} JPM认为，在LNG价格高位、供需偏紧的环境下，中国石油考虑出售LNG Canada部分股权具备合理性，主要逻辑在于回收资本、降低二期资本开支负担。参考Petronas近期以30亿美元出售5\%股权及上游资产，当前估值具有吸引力。中国石油2018年批准项目研究总额34.6亿美元，一期资本开支超预算30\%至260亿加元，二期预计182亿加元，引入合作方可分散风险。市场担忧资产减值，但JPM认为中性偏正面，最新预测已保守计入150亿人民币减值。基本面强劲，预计二季度净利580亿人民币（同比+56\%），上半年1060亿，支撑股价表现。
{\small\color{refgray}数据：LNG Canada一期产能1400万吨/年，2025年Q2调试，2026年下半年满负荷\par}
{\small\color{refgray}数据：二期产能1400万吨/年，资本开支约182亿加元（136.5亿美元）\par}
{\small\color{refgray}数据：Petronas出售5\%股权及上游资产20\%权益，交易额约30亿美元\par}
{\small\color{refgray}数据：中国石油2018年批准项目研究总额34.6亿美元\par}
{\small\color{refgray}边际变化：从单纯持有LNG资产转向利用高景气窗口部分退出，以分担二期资本开支，反映国际油气公司资本配置策略的边际变化。\par}
\subsubsection*{关键数据}
\begin{itemize}
  \item LNG Canada一期产能1400万吨/年，2025年Q2调试，2026年下半年满负荷
  \item 二期资本开支约182亿加元（136.5亿美元），一期超预算30\%至260亿加元
  \item Petronas出售5\%股权及上游资产20\%权益，交易额约30亿美元
  \item 中国石油2018年批准项目研究总额34.6亿美元
  \item 预计二季度净利润580亿人民币，同比+56\%
  \item 上半年净利润约1060亿人民币
  \item 中国石油最新预测已计入150亿人民币减值
\end{itemize}
\begin{figure}[htbp]
\centering
\includegraphics[width=0.88\linewidth]{figures/fig_078.jpg}
\caption{Figure 5 - Figure 5: China monthly gas S-D Figure 6: China GRM vs. SOE/Teapot utilization \% China GRM (Rmb/t) vs SOE and teapot utilization \% \#\# Overweight 0857.HK, 857 HK Price (29 Jul 26):HK\textbackslash{}\$9.93}
\end{figure}
{\small\color{refgray}本节综合 1 篇研究；涉及 JPM。}
\newpage
\subsection{宏观、央行与利率：供给洪峰下的需求韧性、成本冲击与增长分化}
\textbf{全球信用债市场在创纪录供给下利差仍偏紧，需求韧性是核心支撑，但AI相关发行集中度与长久期需求疲软构成利差走阔风险；苹果财报显示需求强劲但供给约束、内存涨价与服务减速同时压制盈利，毛利率压力超预期；被动元件行业景气未变，涨价逻辑延续但短期指引低于预期。}
\subsubsection*{主流共识}
\begin{itemize}
  \item 全球信用债需求端技术面仍能对冲供给不确定性，主要信用市场均录得正净流入，利差维持历史区间偏紧位置。
  \item 苹果iPhone需求依然强劲，六月季度收入创纪录，供给约束是九月季度收入上限的主因而非需求问题。
  \item 被动元件行业AI驱动的供给紧张逻辑未变，涨价假设维持，产能利用率提升与库存下降显示基本面健康。
  \item 内存价格上涨对科技硬件毛利率构成显著压力，且短期难以缓解。
\end{itemize}
\subsubsection*{分歧与条件差异}
\begin{itemize}
  \item 信用债需求结构分歧：被动ETF主导流入而主动产品净流出，长久期品种需求明显落后于短端，AI相关发行集中度可能推动利差走阔。
  \item 苹果九月季度毛利率指引分歧：表面47\%-48\%但剔除关税返还后实际约46.5\%，内存成本贡献超100\%压缩，市场对后续毛利率修复节奏存在分歧。
  \item 被动元件三季度指引分歧：公司'温和增长'指引低于市场预期，但产能利用率跳升与库存下降暗示指引存在上调空间。
  \item 苹果服务业务增速放缓（12\% vs 预期14.6\%）是否趋势性拐点，市场存在分歧。
\end{itemize}
\subsubsection*{机构观点}
\paragraph{GS} 全球信用债需求端技术面仍能对冲创纪录供给，美元IG ETF连续65周净流入，年初至今AUM增约8\%，但结构分化明显：被动ETF主导流入，主动产品净流出，长久期需求疲软（12年以上仅占净需求22\%）。AI相关发行集中度上升，叠加长端供给压力（15年以上占发行27.6\%），利差面临温和走阔风险，且不确定性偏向更大幅度。
{\small\color{refgray}数据：美元IG年初至今发行近1.5万亿美元，同比+31\%\par}
{\small\color{refgray}数据：美元IG ETF连续65周净流入，AUM增约8\%\par}
{\small\color{refgray}数据：3-7年期债券贡献49\%净需求，12年以上仅22\%\par}
{\small\color{refgray}数据：15年以上期限占美国IG发行27.6\%，五年最高\par}
{\small\color{refgray}边际变化：从单纯关注供给冲击转向强调需求结构分化，AI相关发行集中度成为利差走阔的新风险点。\par}
\paragraph{MS} 苹果六月季度iPhone收入同比+22\%至543亿美元创纪录，但九月季度供给约束（3nm）限制收入上限，毛利率指引47\%-48\%含约1个百分点关税返还，剔除后实际约46.5\%，环比降150基点，内存成本贡献超100\%压缩。服务收入同比+12\%低于预期14.6\%，FY27 EPS预期从10.39降至10.00美元，目标价从364降至360美元。
{\small\color{refgray}数据：iPhone六月季度收入543亿美元，同比+22\%\par}
{\small\color{refgray}数据：九月季度毛利率指引47\%-48\%，剔除关税返还后约46.5\%\par}
{\small\color{refgray}数据：内存成本贡献超过100\%的毛利率压缩\par}
{\small\color{refgray}数据：服务收入同比+12\%，低于预期14.6\%\par}
{\small\color{refgray}边际变化：从关注需求端转向供给约束与成本冲击，内存涨价对毛利率压力超预期，服务业务减速成为新关注点。\par}
\paragraph{JPM} 国巨二季度营收444.56亿新台币超预期9\%，毛利率38.5\%，净利率21.2\%，EPS 4.54元同比+88\%。三季度指引'温和增长'隐含环比约5\%，低于市场预期11-12\%，但产能利用率升至90-95\%（二季度80-85\%），库存天数降至110天，指引存在上调空间。维持2026年涨价5-10\%、2027年10-20\%假设，AI供给紧张逻辑未变。
{\small\color{refgray}数据：二季度营收444.56亿新台币，超预期9\%\par}
{\small\color{refgray}数据：三季度指引环比约+5\%，市场预期11-12\%\par}
{\small\color{refgray}数据：产能利用率升至90-95\%，二季度80-85\%\par}
{\small\color{refgray}数据：库存天数从120天降至110天\par}
{\small\color{refgray}边际变化：从关注涨价预期转向指引与稼动率的组合信号，库存下降与产能利用率跳升暗示需求端未转弱。\par}
\subsubsection*{关键数据}
\begin{itemize}
  \item 美元IG年初至今发行近1.5万亿美元，同比+31\%
  \item 美元IG ETF连续65周净流入，AUM增约8\%
  \item 3-7年期债券贡献49\%净需求，12年以上仅22\%
  \item 苹果iPhone六月季度收入543亿美元，同比+22\%
  \item 苹果九月季度毛利率剔除关税返还后约46.5\%，内存成本贡献超100\%压缩
  \item 苹果服务收入同比+12\%，低于预期14.6\%
  \item 国巨二季度营收444.56亿新台币，超预期9\%，EPS同比+88\%
  \item 国巨三季度指引环比约+5\%，市场预期11-12\%
  \item 国巨产能利用率升至90-95\%，库存天数降至110天
\end{itemize}
\begin{figure}[htbp]
\centering
\includegraphics[width=0.88\linewidth]{figures/fig_167.jpg}
\caption{Exhibit 4 - Exhibit 4: Demand for long-end corporate IG bonds has been more tepid than for front-end and intermediate durations Cumulative year-to-date net flows into USD IG funds as a percentage of AUM, by tenor While demand from yield-base}
\end{figure}
\begin{figure}[htbp]
\centering
\includegraphics[width=0.88\linewidth]{figures/fig_168.jpg}
\caption{Exhibit 5 - Exhibit 5: Trade-level data shows a record pace of net buying in USD IG corporate bonds, with USD HY approaching record levels Net customer purchases of cash bonds from dealers by year, for IG (left panel) and HY (right panel)}
\end{figure}
\begin{figure}[htbp]
\centering
\includegraphics[width=0.88\linewidth]{figures/fig_165.jpg}
\caption{Exhibit 1 - Exhibit 2: European credit funds have experienced modest net inflows year to date Year-to-date cumulative net flows into EUR credit funds}
\end{figure}
\begin{figure}[htbp]
\centering
\includegraphics[width=0.88\linewidth]{figures/fig_166.jpg}
\caption{Exhibit 2 - Exhibit 2: European credit funds have experienced modest net inflows year to date Year-to-date cumulative net flows into EUR credit funds Another important nuance is the recent bifurcation of IG flows across the curve. Exhibit 4}
\end{figure}
{\small\color{refgray}本节综合 4 篇研究；涉及 JPM、MS、GS。}
\newpage
\subsection{中国消费修复：底部企稳，结构分化}
\textbf{中国消费活动在6月呈现边际改善，但整体仍低于历史均值，修复集中在必需消费，可选消费与出行链偏弱，大规模刺激政策预期较低。}
\subsubsection*{主流共识}
\begin{itemize}
  \item 消费数据边际改善，但整体修复力度有限，仍低于历史均值。
  \item 商品零售与餐饮等必需消费率先企稳，修复顺序为‘日常消费先稳、可选消费后置’。
  \item 短期内大规模消费刺激政策出台的可能性较低，消费大概率延续温和修复。
\end{itemize}
\subsubsection*{分歧与条件差异}
\begin{itemize}
  \item 消费修复的可持续性存在分歧：部分观点认为底部企稳信号明确，另一部分担忧可选消费（如航空）走弱拖累整体复苏。
  \item 对政策刺激的预期不同：有观点认为基数效应改善将自然推动消费回升，另有观点认为需要额外政策支持才能突破当前温和状态。
\end{itemize}
\subsubsection*{机构观点}
\paragraph{MS} MS的五因子消费活动追踪指标显示，6月Z-score较5月边际改善，但整体仍处于历史均值下方。改善主要由商品零售（不含汽车）和餐饮收入同比回升驱动，而航空客运量三个月移动平均同比回落，表明修复集中在必需性、即时性消费，可选消费与跨区域出行需求仍偏弱。住户贷款与乘用车零售保持平稳，为指数提供底部支撑。MS认为，尽管基数效应逐步改善，消费大概率维持温和修复，短期内大规模刺激政策可能性低。
{\small\color{refgray}数据：五因子Z-score指数6月边际改善，但仍低于历史均值\par}
{\small\color{refgray}数据：商品零售（不含汽车）同比增速回升\par}
{\small\color{refgray}数据：餐饮收入同比好转\par}
{\small\color{refgray}数据：航空客运量三个月移动平均同比回落\par}
{\small\color{refgray}边际变化：相比5月，6月消费活动指数小幅回升，但结构分化加剧：必需消费改善，可选消费（航空）走弱，显示修复力度有限。\par}
\subsubsection*{关键数据}
\begin{itemize}
  \item 6月五因子消费活动Z-score边际改善，但仍低于历史均值
  \item 商品零售（不含汽车）同比增速回升
  \item 餐饮收入同比好转
  \item 航空客运量三个月移动平均同比回落
  \item 住户贷款同比平稳
  \item 乘用车零售三个月移动平均平稳
  \item 大规模消费刺激政策可能性较低
\end{itemize}
\begin{figure}[htbp]
\centering
\includegraphics[width=0.88\linewidth]{figures/fig_087.jpg}
\caption{Exhibit 1 - Exhibit 1: China Five-factor Consumer Activity Z-Score vs. MSCI China YoY change – Consumer activity improved marginally in June but remained weak MS ASIA LIMITED+ Laura Wang Equity Strategist}
\end{figure}
{\small\color{refgray}本节综合 1 篇研究；涉及 MS。}
\newpage
\subsection{电气化与AI基建：执行力拐点与数据中心结构性增长}
\textbf{施耐德电气上半年业绩证明其执行力拐点已至，数据中心订单三位数增长叠加原材料通胀回落，推动有机增长全面超预期，公司正从执行调整期转向改善期，并在预制化AI基础设施领域构建隐形壁垒。}
\subsubsection*{主流共识}
\begin{itemize}
  \item 数据中心相关业务是电气设备板块结构性增长的核心驱动力，订单增速显著高于传统业务。
  \item 原材料（铜、银）通胀压力在下半年将环比回落，利好毛利率扩张。
  \item 施耐德电气上半年有机增长全面超预期，各地区各业务均实现正增长，显示增长质量高。
  \item 定价弹性是短期利润率改善的关键，H2定价有望继续加速。
\end{itemize}
\subsubsection*{分歧与条件差异}
\begin{itemize}
  \item 市场对施耐德电气经营杠杆的持续性存疑，但报告认为执行力改善已获数据验证，分歧在于改善能否延续至2026年。
  \item 关于原材料套保比例假设不同（报告假设20\%），若实际套保更高或更低，H2通胀回落幅度存在差异。
  \item 数据中心增长是否可持续：部分观点担忧AI资本开支周期性，但报告强调预制化方案和合作伙伴关系（如富士康）提供长期壁垒。
  \item 对施耐德电气估值溢价的分歧：目标价上调至330欧元，但市场对估值消化能力看法不一。
\end{itemize}
\subsubsection*{机构观点}
\paragraph{Bernstein} 施耐德电气上半年业绩验证执行力拐点，LFL销售增长16.5\%超预期6个百分点，能源管理北美增长24.8\%，数据中心订单三位数增长。定价弹性加速（Q2为Q1的2.5倍），H2原材料通胀环比回落，毛利率全年扩张120个基点。预制化方案与富士康合作构建AI基础设施壁垒，上调2026年有机增长预期280个基点，目标价310升至330欧元。
{\small\color{refgray}数据：LFL销售增长16.5\%，超预期约6个百分点\par}
{\small\color{refgray}数据：能源管理北美LFL增长24.8\%，超预期840个基点\par}
{\small\color{refgray}数据：数据中心订单二季度三位数增长\par}
{\small\color{refgray}数据：H1定价约2.8\%，Q2定价为Q1的2.5倍\par}
{\small\color{refgray}边际变化：从过去12个月偏弱经营杠杆的审视，转向执行力改善的确认，且数据中心增长与原材料回落形成双重支撑，预期上调。\par}
\subsubsection*{关键数据}
\begin{itemize}
  \item 施耐德电气LFL销售增长16.5\%，超预期约6个百分点
  \item 能源管理北美LFL增长24.8\%，超预期840个基点
  \item 工业自动化中国与东亚LFL增长19.6\%，超预期12个百分点
  \item 数据中心订单二季度三位数增长
  \item H1定价约2.8\%，Q2定价为Q1的2.5倍
  \item 全年毛利率扩张120个基点（剔除关税退税后80个基点）
  \item 2026年有机增长预期上调280个基点
  \item 目标价从310欧元上调至330欧元
\end{itemize}
\begin{figure}[htbp]
\centering
\includegraphics[width=0.88\linewidth]{figures/fig_112.jpg}
\caption{EXHIBIT 2 - EXHIBIT 3: We estimate Schneider's raw material exposure using reports and peer benchmarking. Layering on expected H2'26 inflation (at spot prices) implies a sequential step-down in raw material inflation. We assume Schneider hedge}
\end{figure}
\begin{figure}[htbp]
\centering
\includegraphics[width=0.88\linewidth]{figures/fig_113.jpg}
\caption{EXHIBIT 6 - EXHIBIT 6: Schneider's 800 VDC portfolio. \#\# 800 VDC portfolio designed in step with the next generation of AI Infrastructure \#\# 800 VDC Power Rack (Sidecar) \textbackslash{}- Location: IT Rack • Tech: Power Shelves • Power: Up to 1.2 MW DC to}
\end{figure}
\begin{figure}[htbp]
\centering
\includegraphics[width=0.88\linewidth]{figures/fig_114.jpg}
\caption{EXHIBIT 8 - EXHIBIT 8: Foxconn strategic partnership to accelerate modular data center offering. \#\# Accelerating AI infrastructure deployment with Foxconn strategic partnership}
\end{figure}
\begin{figure}[htbp]
\centering
\includegraphics[width=0.88\linewidth]{figures/fig_115.jpg}
\caption{EXHIBIT 8 - EXHIBIT 8: Foxconn strategic partnership to accelerate modular data center offering. \#\# Accelerating AI infrastructure deployment with Foxconn strategic partnership Foxconn's manufacturing power • Foxconn brings advanced compute,}
\end{figure}
{\small\color{refgray}本节综合 1 篇研究；涉及 Bernstein。}
\newpage
\subsection{新兴市场增长空间分化，加拿大鹅收入质量待观察}
\textbf{新兴市场经济体产出缺口分化明显，多数仍有闲置产能，半导体驱动的增长尚未引发广泛过热；加拿大鹅单季收入超预期但质量打折，DTC客流是后续关键。}
\subsubsection*{主流共识}
\begin{itemize}
  \item 多数新兴市场经济体仍在潜在产出水平附近或以下运行，存在闲置产能。
  \item 半导体带动的增长集中在少数企业，未传导至整体劳动力市场，不构成广泛宏观过热。
  \item 工资增长与产出缺口的相关性依然较强，但产出缺口与通胀的关系因全球供应链影响而弱化。
\end{itemize}
\subsubsection*{分歧与条件差异}
\begin{itemize}
  \item 区域分化：CEEMEA和中国产出缺口进一步走负，拉美接近潜力，亚洲（除中国）已转正。
  \item 方法论差异：传统滤波方法可能高估台湾等地的正缺口，生产函数估算显示增长外溢有限。
  \item 加拿大鹅收入超预期主要来自批发时间性回补，而非DTC内生需求，市场对收入质量存在分歧。
\end{itemize}
\subsubsection*{机构观点}
\paragraph{GS} GS认为加拿大鹅F1Q27收入超预期主要源于批发订单时间前移和亚太动能，但DTC同店销售同比下滑3.2\%（上季+10\%），客流疲软尤其在欧洲和美国，收入质量打折。关税若落地可能影响FY27利润率约200bps。维持目标价C\$11.50，核心关注秋冬旺季门店客流能否回暖。
{\small\color{refgray}数据：收入\$118.9mn，同比+10.3\%，超预期（预期-1.9\%）\par}
{\small\color{refgray}数据：批发收入+66.5\%至\$29.8mn，部分为订单前移\par}
{\small\color{refgray}数据：DTC同店销售-3.2\%，上季+10.0\%\par}
{\small\color{refgray}数据：调整后EPS -C\$0.89，好于预期-C\$0.95\par}
{\small\color{refgray}边际变化：从单纯关注收入超预期转向强调收入质量，DTC同店由正转负是边际变化。\par}
\paragraph{GS} GS更新22个新兴市场产出缺口估算，显示多数经济体仍有闲置产能，但区域分化明显：CEEMEA和中国缺口走负，拉美接近潜力，亚洲（除中国）转正。半导体驱动的增长集中在少数企业，未传导至劳动力市场，韩国缺口仍为负，台湾、香港、马来西亚仅略正。工资增长与缺口相关性仍强，但通胀受全球供应链影响更大。
{\small\color{refgray}数据：覆盖22个新兴市场\par}
{\small\color{refgray}数据：韩国产出缺口为负\par}
{\small\color{refgray}数据：台湾、香港、马来西亚缺口略正\par}
{\small\color{refgray}数据：中国存在较大闲置产能\par}
{\small\color{refgray}边际变化：相比传统滤波方法，生产函数估算显示半导体增长外溢有限，避免高估过热风险。\par}
\subsubsection*{关键数据}
\begin{itemize}
  \item 加拿大鹅收入\$118.9mn，同比+10.3\%，超预期
  \item DTC同店销售-3.2\%，上季+10.0\%
  \item 批发收入+66.5\%，部分为订单前移
  \item 关税影响FY27利润率约200bps
  \item 22个新兴市场产出缺口估算
  \item 韩国产出缺口为负，台湾、香港、马来西亚略正
  \item 中国闲置产能显著
  \item 巴西产出高于潜力，阿根廷、墨西哥、智利有闲置
\end{itemize}
\begin{figure}[htbp]
\centering
\includegraphics[width=0.88\linewidth]{figures/fig_151.jpg}
\caption{Exhibit 1 - Exhibit 1: Our Regional Output Gap Estimates Are Either Close to Potential (EM Asia ex China and LatAm) or Below Potential (CEEMEA and China) Exhibit 2 compares the latest (2026Q2) readings. Our results suggest significant spare}
\end{figure}
\begin{figure}[htbp]
\centering
\includegraphics[width=0.88\linewidth]{figures/fig_152.jpg}
\caption{Exhibit 2 - Exhibit 2 compares the latest (2026Q2) readings. Our results suggest significant spare capacity in South Africa, CEE ex Poland, South Korea, China, Chile and Argentina. Since our last update, Taiwan's output gap has turned from slightly negative to slightly po}
\end{figure}
\begin{figure}[htbp]
\centering
\includegraphics[width=0.88\linewidth]{figures/fig_153.jpg}
\caption{Exhibit 3 - Exhibit 3: Wage Growth is Generally Higher in Economies that are Running Positive Output Gaps Core inflation (wage) gaps are defined as deviations of core inflation (wage growth) from target (target-consistent levels) and are based}
\end{figure}
\begin{figure}[htbp]
\centering
\includegraphics[width=0.88\linewidth]{figures/fig_154.jpg}
\caption{Exhibit 3 - Exhibit 3: Wage Growth is Generally Higher in Economies that are Running Positive Output Gaps Core inflation (wage) gaps are defined as deviations of core inflation (wage growth) from target (target-consistent levels) and are based}
\end{figure}
{\small\color{refgray}本节综合 2 篇研究；涉及 GS。}
\newpage
\subsection{中国汽车互联网平台：投入加码与资本回报的再平衡}
\textbf{2026年中报窗口，中国汽车互联网平台进入投入强度与股东回报并重的阶段：途虎养车加速门店扩张但利润率承压，汽车之家以新零售和回购对冲传统业务下滑。行业层面，燃油车销量降幅收窄与新能源渗透率阶段性回调并存，平台需在需求节奏变化中验证投入的ROI。}
\subsubsection*{主流共识}
\begin{itemize}
  \item 燃油车销量同比降幅收窄，年底有望转正，但新能源渗透率从2H25的58\%回落至1H26的54\%，叠加保养周期延长，传统汽车媒体与售后市场空间被压缩。
  \item 两家平台均加大投入（途虎门店扩张、汽车之家新零售），但短期利润被侵蚀，市场关注投入回报率而非单纯增长。
  \item 股东回报成为共同主题：途虎回购不低于5000万股，汽车之家合计回报率约22\%并新增4亿美元回购。
\end{itemize}
\subsubsection*{分歧与条件差异}
\begin{itemize}
  \item 途虎选择重资产门店扩张（上半年净增900家），而汽车之家走轻资产加盟模式（600家店），两种路径的资本效率与风险不同。
  \item 对投入ROI的判断分歧：GS维持中性评级，认为扩张成本前置超预期，但长期增长潜力与股东回报可能平衡风险。
  \item 新能源渗透率回调是暂时性还是趋势性，影响平台收入结构转型的紧迫性。
\end{itemize}
\subsubsection*{机构观点}
\paragraph{GS} GS认为中国汽车互联网平台进入投入与回报的验证期。途虎养车2026年上半年净增约900家门店，门店总数同比+24\%，但毛利率下降和销售费用上升导致调整后净利润预测平均下调17\%，调整后经营利润率从4.1\%降至1.4\%。汽车之家传统业务承压，收入预期下调8\%，但新零售品牌'家家好车'计划开600家店，二手车出口业务2H26放量。两家公司均通过回购提升股东回报，但投入ROI是股价关键。
{\small\color{refgray}数据：途虎2026H1净增约900家门店，同比+24\%\par}
{\small\color{refgray}数据：途虎调整后经营利润率从4.1\%降至1.4\%\par}
{\small\color{refgray}数据：汽车之家收入预期下调8\%\par}
{\small\color{refgray}数据：汽车之家2026年计划开600家店\par}
{\small\color{refgray}边际变化：从单纯追求增长转向强调投入回报率，且股东回报（回购）成为新叙事。\par}
\subsubsection*{关键数据}
\begin{itemize}
  \item 国内成品油价格年内上涨8\%
  \item 新能源渗透率从2H25的58\%回落至1H26的54\%
  \item 途虎2026H1净增约900家门店，同比+24\%
  \item 途虎调整后净利润率从5.2\%降至2.5\%
  \item 途虎调整后经营利润同比下滑61\%
  \item 途虎回购不低于5000万股，占总股本6\%
  \item 途虎账上净现金70亿，相当于市值80\%以上
  \item 汽车之家2026年计划开600家店，主攻下沉市场
  \item 汽车之家2026年合计回报率约22\%，新增4亿美元回购
\end{itemize}
\begin{figure}[htbp]
\centering
\includegraphics[width=0.88\linewidth]{figures/fig_137.jpg}
\caption{Exhibit 2 - Exhibit 4: Tuhu DAU increased +14\% yoy in Jun (vs. +12\% yoy in May), while Tuhu DAU/Tuhu Merchant app DAU reached 26.4 in Jun (vs. 25.0 in May)}
\end{figure}
\begin{figure}[htbp]
\centering
\includegraphics[width=0.88\linewidth]{figures/fig_139.jpg}
\caption{Exhibit 3 - Exhibit 5: We expect Tuhu's LTM transacting user to grow +18\% yoy to c.31.3mn in 1H26E \#\# Autohome (ATHM/2518.HK)}
\end{figure}
\begin{figure}[htbp]
\centering
\includegraphics[width=0.88\linewidth]{figures/fig_138.jpg}
\caption{Exhibit 2 - Exhibit 4: Tuhu DAU increased +14\% yoy in Jun (vs. +12\% yoy in May), while Tuhu DAU/Tuhu Merchant app DAU reached 26.4 in Jun (vs. 25.0 in May) Exhibit 5: We expect Tuhu's LTM transacting user to grow +18\% yoy to c.31.3mn in 1H26}
\end{figure}
\begin{figure}[htbp]
\centering
\includegraphics[width=0.88\linewidth]{figures/fig_140.jpg}
\caption{Exhibit 4 - Exhibit 4: Tuhu DAU increased +14\% yoy in Jun (vs. +12\% yoy in May), while Tuhu DAU/Tuhu Merchant app DAU reached 26.4 in Jun (vs. 25.0 in May) Exhibit 5: We expect Tuhu's LTM transacting user to grow +18\% yoy to c.31.3mn in 1H26}
\end{figure}
{\small\color{refgray}本节综合 1 篇研究；涉及 GS。}
\newpage
\section{辅助信号｜战略咨询}
本板块聚焦波士顿咨询近期关于并购周期、AI投入与组织能力的研究，提炼对宏观、产业及企业战略具有边际意义的信号。核心发现包括：弱宏观环境下并购的长期回报显著优于强周期；AI投入的保守正从风险管理转为竞争劣势；领导团队的信念差距是制约组织业绩的隐性变量。
\subsection{弱周期并购的长期回报优势与执行框架}
\textbf{经济下行期并非并购的禁区，反而因资产估值低、竞争减少而提供更高的长期回报；但成功取决于交易类型、组织能力和执行纪律，而非宏观环境本身。}
\begin{itemize}
  \item 波士顿咨询研究显示，弱宏观环境周期完成的并购交易，两年后收购方相对股东总回报（RTSR）比强周期高出超过9个百分点，差距在交割后第一年即显现并持续扩大。
  \item 2018年收购方公告日累计异常回报（CAR）均值转负（-0.4\%），为2012年以来首次，表明市场短期定价与中期价值创造存在系统性偏差，决策者应区分CAR与RTSR的不同含义。
  \item 弱周期中非核心（跨界）收购的一年期回报比核心收购高出3.9个百分点，反映竞争性买家减少和估值折价带来的价值空间。
  \item 2020年疫情调查显示，73\%的受访者经历并购估值重新谈判，但到夏末并购活动反弹快于预期；报告建议买方在早期阶段重做尽职调查、测试情景敏感度（V/U/L型复苏），卖方则需结合内部与外部视角判断交易是否重启、调整或终止。
  \item 报告强调，并购回报由可重复的流程决定，而非单笔交易；企业应建立持续的目标筛选机制，而非在周期底部犹豫或条件反射式终止计划。
\end{itemize}
\begin{figure}[htbp]
\centering
\includegraphics[width=0.88\linewidth]{figures/fig_285.jpg}
\caption{EXHIBIT 1 - The rapid rise in remote work and virtual schooling is another COVID-19 effect stimulating technology investment opportunities. In the education sector, government action to keep K-12 schooling happening at home is significantly accelerating edtech digital inn}
\end{figure}
\subsection{AI投入的竞争门槛与组织变革价值}
\textbf{AI投入的保守正从风险管理转为竞争劣势；深度整合AI的企业在成本、利润和资本回报上形成显著优势，且价值创造主要来自组织与流程变革而非算法本身。}
\begin{itemize}
  \item 波士顿咨询2026年7月报告显示，AI投入占营收比例较高的企业（约1.7\%）成本降幅达同行三倍，利润率高出60\%，投入资本回报率是同行2.7倍，呈现明显门槛效应。
  \item AI价值构成中约70\%来自组织与流程变革，仅10\%来自算法；企业需用短周期验证机制替代静态商业案例评估，以动态能力建设为核心。
  \item 报告警告，继续聚焦AI成本控制、仅运行低成本试点的企业将逐渐落后；优势来自将AI嵌入核心业务流程并重新设计运营模式，具有复利特征。
\end{itemize}
\subsection{领导团队的信念差距与组织业绩上限}
\textbf{团队业绩不佳常被归因于外部因素，但波士顿咨询调查显示，领导团队共享的未经检验假设才是限制组织上限的隐性变量；重构信念是回报最高的领导力杠杆。}
\begin{itemize}
  \item 2026年3月对11个市场563名高级领导者的调查显示，63\%的受访者认为所在组织的高级领导者并未为长期成功做好准备；同时做到重构信念、刻意练习、固化行为的团队，业绩领先同行可能性高出144\%。
  \item 重构信念是三项能力中回报最高的一环：主动重构信念的领导者业绩领先可能性高出168\%，且信念只能通过体验建立，无法靠逻辑说服。
  \item 团队共享的假设比个人盲点更难移动，单一领导者改变后若回到未变的团队，旧行为会将其拉回；百思买前CEO休伯特·乔利的转型案例印证了先改信念再改行为的路径。
\end{itemize}
\begin{figure}[htbp]
\centering
\includegraphics[width=0.88\linewidth]{figures/fig_282.jpg}
\caption{Exhibit 1 - Most leadership teams carry these kinds of unexamined assumptions, whether about the market, their own limits, or what good leadership demands of them. In a recent BCG survey of 563 senior leaders across 11 markets, \$63\textbackslash{}\%\$ said that the senior leaders in their}
\end{figure}
\begin{figure}[htbp]
\centering
\includegraphics[width=0.88\linewidth]{figures/fig_283.jpg}
\caption{Exhibit 2 - Consider the two things that leaders most often credit for their own growth. Having a growth mindset (63\%) and embracing discomfort and uncertainty (51\%) may sound like capabilities, but both reflect assumptions that leaders hold about themselves. (See Exhibit}
\end{figure}
\newpage
\section{辅助信号｜智库/国际机构}
本板块聚焦布鲁盖尔研究所、IMF、经合组织与世贸组织的最新研究，提炼出欧洲金融监管与银行业重构、全球宏观与债务治理、以及资源型经济与结构性转型三大主题，为宏观与资产配置提供辅助信号。
\subsection{欧洲金融监管与银行业重构}
\textbf{欧盟金融监管架构正面临碎片化与不确定性挑战，银行处置清算规则差异、监管职能分离及欧元区治理改革方向成为核心议题。}
\begin{itemize}
  \item 布鲁盖尔研究所建议欧盟以“双峰监管”为长期愿景，分离审慎监管与行为监管；截至2015年，欧盟31\%的银行资产和36\%的保险资产归属于金融集团，现行按行业划分的监管体系难以应对混业经营。
  \item 欧盟银行处置与清算适用不同法律框架：处置程序要求至少8\%的内部纾困，而清算仅需“轻量”负担分担；意大利两银行清算与西班牙大众银行处置的对比显示，公共资金介入规则差异显著。
  \item 欧元区财政治理改革存在两派方案：容克主张强化欧盟委员会权力，朔伊布勒提议升级欧洲稳定机制；布鲁盖尔提出中间方案，将欧元集团改造为“财政政策欧元体系”，以常设主席制强化整体利益代表性。
  \item 欧洲央行货币政策正常化将进入“新常态”：资产负债表规模较危机前翻四倍，被动持有到期需约14年才能回到危机前水平，中性利率下降意味着未来更依赖资产负债表工具而非利率调整。
  \item 欧洲银行业不良贷款存量约8700亿欧元，潜在供应接近2万亿欧元，但2016年二级市场交易额仅1460亿欧元；估值分歧、投资者集中与信息不对称构成市场失灵的三重障碍。
\end{itemize}
\begin{figure}[htbp]
\centering
\includegraphics[width=0.88\linewidth]{figures/fig_192.jpg}
\caption{Figure 1 - Figure 1: Supervisory synergies and conflicts The first potential conflict of interest between systemic supervision and prudential supervision relates to the possibility of lender-of-last-resort operations (LoLR) by the central}
\end{figure}
\begin{figure}[htbp]
\centering
\includegraphics[width=0.88\linewidth]{figures/fig_194.jpg}
\caption{Figure 1 - Figure 1: Use of public funds in the EU framework If a bank is declared by the ECB to be failing or likely to fail \$\textasciicircum{}\{11\}\$ , the precautionary recapitalisation option is not available, and the choice is between liquidation or re}
\end{figure}
\subsection{全球宏观与债务治理}
\textbf{全球增长预测存在系统性乐观偏差，主权债务偿还存在事实优先级，金融压抑可能回归，这些结构性因素重塑宏观与债务风险定价。}
\begin{itemize}
  \item IMF对1999-2023年WEO预测的检验显示，增长预测平均高出实际约0.5个百分点，且美国与中国增长溢出被显著低估：实际同步弹性为0.84，预测中仅0.34，预测碎片化问题突出。
  \item IMF基于119个新兴市场和发展中经济体1980-2022年数据发现，主权债务偿还存在事实优先级：IMF和世界银行最优先，官方双边债权人次之，债券持有人和商业银行居后，其中商业银行优先级最低。
  \item IMF研究显示，金融压抑并非历史遗物：二战后达到顶峰，资本账户自由化时期消退，2008年全球金融危机后重新上升；银行与家庭持有政府债务的缺口可作为量化压抑强度的指标。
  \item IMF对英国劳动税的分析指出，统一提高劳动税率空间有限，边际有效税率在收入阈值处存在陡峭跳跃，高收入者边际税率高于其他G7国家，税制高度累进性集中于收入分布上端。
  \item IMF研究揭示，新加坡私人消费占GNDI比重从1990年的46\%降至2024年的38\%，在31个发达经济体中最低；老年抚养比、劳动收入份额、金融中心强度等结构性因素而非周期性因素主导了这一趋势。
\end{itemize}
\begin{figure}[htbp]
\centering
\includegraphics[width=0.88\linewidth]{figures/fig_259.jpg}
\caption{Figure 1 - Figure 1: Stability of medium-term anchors over time. The grey dots are individual country medium-term forecasts (five years ahead, \$v+5\$ ) and the shaded band their interquartile range; the solid blue line is the global median an}
\end{figure}
\begin{figure}[htbp]
\centering
\includegraphics[width=0.88\linewidth]{figures/fig_229.jpg}
\caption{Figure 1 - Figure 1: Scatter Plots of Sovereign Bond Spreads and Debt Composition by Creditor Type}
\end{figure}
\subsection{资源型经济与结构性转型}
\textbf{大宗商品繁荣、资源依赖与制度落差共同塑造新兴市场的增长路径，AI与可持续倡议则成为发达经济体生产率与治理变革的新变量。}
\begin{itemize}
  \item IMF基于智利2003-2013年微观数据发现，铜价大涨期间矿业全要素生产率下降约8\%，资源错配可解释其中约一半损失；出口暴露度高的企业收入增长但生产率未相应提升，低生产率企业扩张更快。
  \item IMF对毛里塔尼亚的技术援助显示，针对资源依赖型国家定制的债务动态工具（DDT RRC）将非采掘初级余额作为核心财政变量，并纳入自然灾害模块，使债务预测能区分资源收入波动与政策调整的贡献。
  \item IMF报告指出，英国劳动力中近三分之二从事AI可暴露职业，且超过一半属于高互补性岗位，AI中期TFP年增益可达0.1-0.5个百分点；但监管不确定性和AI技能短缺是最紧迫的约束，若减半监管约束并扩大技能供给，产出增益可提高三分之二。
  \item OECD对1078项可持续倡议的分析显示，仅30\%存在政府结构性参与，71\%被政府在非约束性政策文件中提及，呈现“轻介入、重引用”的关系，可持续倡议扩张主要由市场力量驱动。
  \item 布鲁盖尔研究所指出，欧盟煤电仍占发电量约四分之一，却贡献电力部门75\%的二氧化碳排放；通过改造现有全球化调整基金每年投入1.5亿欧元支持产煤地区转型，可绕开条约限制推动煤电退出。
\end{itemize}
\begin{figure}[htbp]
\centering
\includegraphics[width=0.88\linewidth]{figures/fig_219.jpg}
\caption{Figure 1 - Figure 1: Global Copper Price and Chile's Mining TFP Our identification exploits differential price changes across commodity products, driven by China's WTO accession and rapid industrialization (Fernández et al., 2023). We assi}
\end{figure}
\begin{figure}[htbp]
\centering
\includegraphics[width=0.88\linewidth]{figures/fig_249.jpg}
\caption{Figure 3 - Figure 3: Contributors to Changes in Public Debt Note: The projections presented in this figure are for illustrative purposes only. \#\# C. Natural Disasters Public Debt Dynamics Tools (ND DDT) \#\# 20. The mission also presented th}
\end{figure}
\subsection{宏观政策与金融稳定}
\textbf{用于校准利率、通胀、财政约束和金融稳定叙事。}
\begin{itemize}
  \item 欧元区扰动爆发十余年后，关于“扰动为何发生、责任在谁”的公共叙事仍然没有收敛。布鲁盖尔研究所发布的一份政策研究显示，德国、法国、意大利和西班牙四国的主流报纸在归因上各说各话，几乎没有形成跨国的共同理解。这份研究基于2007年至2016年间四国四家代表性报纸的51,714篇宏观环境...
  \item 2018年2月，欧盟委员会为塞尔维亚和黑山设定了2025年的入盟参考时间表。这份由布鲁盖尔研究所发布的报告，把视线拉回西巴尔干地区过去二十年的宏观环境轨迹：收入在追赶，但速度在节奏变化；失业率在下降，但年轻人仍被排除在劳动力市场之外；制度在靠近欧盟标准，但法治和反腐败的进展远未达...
  \item 过去三十年，跨国收入差距收窄被普遍视为后发地区追赶成功的证据。IMF 宏观环境学家 Patrick A.
  \item 美元在国际贸易计价和全球资金负债中的双重主导地位，通常被分开研究。贸易计价影响汇率传导与支出转换，负债美元化影响资产负债表与资金脆弱性。这份IMF工作论文提出的核心判断是：两者的交互作用，才是决定汇率不确定性属性、货币不确定性溢价乃至小型开放地区通胀动态的关键机制。
  \item 汇率变动究竟在多大程度上反映宏观环境基本面，又在多大程度上源于与基本面无关的资金摩擦？这个问题长期困扰开放宏观宏观环境学。IMF在2026年7月发布的工作论文中，提出了一种结合符号约束与叙事约束的识别方法，试图从利率平价偏离中分离出纯粹的资金冲击。基于巴西和智利两国近二十年的月度...
  \item 2026年7月发布的IMF能力发展研究所技术援助报告显示，毛里塔尼亚公共债务委员会已具备独立编制债务预测与情景分析的能力。这项为期一年的项目由日本政府资助，于2024年1月至2025年1月间执行，核心成果是将IMF的公共债务动态工具（DDT）本地化，并针对该国资源依赖型宏观环境结...
  \item 2022年以来，英国国债回报表现——尤其是长端——持续抬升并超过G7其他地区。IMF在2026年7月发布的国别报告中拆解了这一现象的构成：长期回报表现上升并非单纯由货币政策预期推动，期限溢价——即超出短期利率预期路径的不确定性补偿——才是主要变量。报告基于2022年9月英国国债市...
\end{itemize}
\begin{figure}[htbp]
\centering
\includegraphics[width=0.88\linewidth]{figures/fig_209.jpg}
\caption{Figure 1 - Figure 2: Le Monde, frequency of selected topics}
\end{figure}
\begin{figure}[htbp]
\centering
\includegraphics[width=0.88\linewidth]{figures/fig_214.jpg}
\caption{Figure 1 - Figure 1: GDP per capita in current international \textbackslash{}\$, PPP adjusted, Germany = 100\%, 2000-16 Figure 2: Real GDP growth, annual percent change, 2000-16 The income convergence process was particularly strong between 2000 and 2009}
\end{figure}
\subsection{气候、发展与社会结构}
\textbf{用于观察长期增长、资源约束、社会结构和政策执行质量。}
\begin{itemize}
  \item 乌兹别克斯坦在2026年完成加入世界贸易组织进程的目标，正在从时间表变成可核验的谈判进度。该国政府代表在最新一次工作组会议上给出的信息是：书面问题数量持续下降、双边市场准入协议已签署31份、与WTO规则对齐的法律文件超过190项。这三个数字放在一起，指向一个明确判断——入世谈判已...
\end{itemize}
\section{结语}
后续需验证的数据与叙事节点包括：微软等超大规模平台下季度资本开支指引是否延续扩张；存储供需缺口在2027年是否如GS预期大于2026年；小米SkyNomad实际交付量能否达到乐观情形；中国8月出口与PMI数据能否验证内外需分化收敛；美联储8月会议对曲线扭曲陡峭化的表态；以及电力设备订单能否持续兑现至2030年代。
\newpage
\section{覆盖概览}
投行/券商 92 篇；战略咨询 5 篇；智库/国际机构 23 篇。\par
投行覆盖：GS 29篇 / JPM 21篇 / Citi 10篇 / MS 10篇 / Bernstein 8篇 / NOM 7篇 / BARC 4篇 / UBS 3篇\par
\section{Disclaimer}
{\scriptsize\color{refgray}
This community is a paid learning and information-sharing community created for general educational purposes only. The membership fee is charged solely for access to community discussions, educational content, organization, and operational costs. It is not an advisory fee, management fee, performance fee, brokerage fee, or compensation for personalized financial, investment, legal, tax, or professional advice.\par
I participate and share information solely as an individual and for educational discussion among community members. I am not acting as a financial adviser, investment adviser, broker, dealer, portfolio manager, legal adviser, tax adviser, accountant, fiduciary, or any other licensed professional adviser. Nothing shared in this community constitutes, and should not be interpreted as, financial advice, investment advice, legal advice, tax advice, a recommendation, solicitation, offer, endorsement, instruction, or guarantee to buy, sell, hold, short, trade, or invest in any security, cryptocurrency, fund, derivative, strategy, or other financial product.\par
All content shared in this community, including messages, discussions, links, files, charts, examples, market commentary, watchlists, AI-generated or AI-polished summaries, and third-party materials, is general, impersonal, and educational in nature. It does not take into account any individual's financial situation, investment objectives, risk tolerance, investment horizon, liquidity needs, tax status, legal restrictions, or suitability requirements. No content should be relied upon as suitable for any specific person, account, portfolio, or financial decision.\par
Information shared in this community may be incomplete, inaccurate, outdated, speculative, AI-generated, AI-edited, or based on third-party internet sources that have not been independently verified. I make no representation or warranty as to the accuracy, completeness, timeliness, reliability, or suitability of any information shared.\par
Financial markets involve substantial risk, including the possible loss of principal. Past performance, historical data, hypothetical examples, simulated results, backtests, personal opinions, or educational examples do not guarantee future results. Each member is solely responsible for conducting their own research, exercising independent judgment, and making their own decisions. Before making any financial, investment, legal, or tax decision, members should consult qualified and licensed professionals.\par
Participation in this community, payment of any membership fee, or communication with me or other members does not create any adviser-client, fiduciary, professional, contractual, agency, partnership, employment, or other advisory relationship. By joining or participating in this community, you acknowledge and agree that you are solely responsible for your own decisions, actions, transactions, profits, losses, risks, and consequences, and that I shall not be liable for any direct, indirect, incidental, consequential, financial, legal, tax, or other loss or damage arising from or related to any content, discussion, information, or material shared in this community.\par
}
\newpage
\begin{center}
{\Large 更多详情报告kcdesk.com}\par\vspace{0.8cm}
\includegraphics[width=0.58\linewidth]{\detokenize{../../prompts/zsxq_img.jpg}}
\end{center}
\end{document}
